\documentclass[11pt,a4paper]{article}
\usepackage[utf8]{inputenc}
% lmodern must precede fontenc: without a scalable Type1 font the T1 encoding
% falls back to bitmap fonts, and microtype's expansion then aborts the run
% with "auto expansion is only possible with scalable fonts".
\usepackage{lmodern}
\usepackage[T1]{fontenc}
\usepackage[margin=1in]{geometry}
\usepackage{amsmath,amssymb}
\usepackage{graphicx}
\usepackage{microtype}
\usepackage[hidelinks,breaklinks]{hyperref}
\usepackage{url}
\urlstyle{same}

% Generated by code/tools/md_to_tex.py from the SurveyForge Markdown output.
% Citation numbers match the source .md exactly.

\title{A Comprehensive Survey on Alignment of Diffusion Models}
\author{Generated by SurveyForge}
\date{}

\begin{document}
\maketitle
\tableofcontents
\newpage



\section{Introduction}


The concept of diffusion models has garnered significant attention in recent years, particularly in the realm of generative modeling. As noted in \cite{doi:10.1145/3626235}, diffusion models represent a class of deep generative models that have achieved state-of-the-art performance in various applications, including image synthesis, video generation, and molecule design. At their core, diffusion models involve a forward process that progressively adds noise to the input data and a reverse process that learns to denoise the data, effectively modeling the underlying distribution of the data.

The importance of alignment in diffusion models cannot be overstated, as it directly impacts the quality and realism of the generated samples. Alignment refers to the degree to which the generated samples conform to the desired characteristics, such as text prompts or user preferences. As highlighted in \cite{arxiv:1809.01036}, achieving robust alignment is crucial for ensuring that the generated samples meet the required standards. However, this is often a challenging task, particularly when dealing with complex data distributions or limited training data.

Recent advances in diffusion models have led to significant improvements in alignment, with techniques such as classifier-free guidance and reinforcement learning from human feedback showing great promise. For instance, \cite{arxiv:2406.08070} introduces a novel approach to classifier-free guidance that tackles the off-manifold challenges inherent in traditional classifier-free guidance, resulting in improved sample quality and invertibility. Similarly, \cite{arxiv:2305.13301} demonstrates the effectiveness of reinforcement learning in fine-tuning diffusion models to optimize downstream reward functions.

Despite these advances, there are still several challenges associated with aligning diffusion models, including mode collapse, training instability, and the need for more efficient optimization methods. As noted in \cite{arxiv:2502.06805}, the high computational cost of training diffusion models limits their practical applications, highlighting the need for more efficient techniques. Furthermore, \cite{arxiv:2303.07576} emphasizes the importance of developing more effective methods for aligning diffusion models with human preferences, particularly in the context of natural language processing.

In addition to these challenges, there are also several emerging trends and future directions in the field of diffusion models, including the integration of diffusion models with other generative models, such as GANs and VAEs, and the application of diffusion models to new domains, such as video generation and 3D modeling. As highlighted in \cite{doi:10.1109/iccv51070.2023.00713}, the development of more versatile and flexible diffusion models that can handle multiple tasks and domains is an exciting area of research. Moreover, \cite{arxiv:2305.13655} demonstrates the potential of combining diffusion models with large language models to enhance prompt understanding and generation quality.

In conclusion, the field of diffusion models is rapidly evolving, with significant advances being made in alignment, efficiency, and applications. As researchers continue to push the boundaries of what is possible with diffusion models, it is essential to address the challenges and limitations associated with these models, including mode collapse, training instability, and the need for more efficient optimization methods. By synthesizing the existing literature and identifying emerging trends and future directions, this survey aims to provide a comprehensive overview of the current state of diffusion models and their applications, as well as insights into the potential future developments in this exciting field. As noted in \cite{doi:10.1109/tpami.2025.3545047}, the potential of diffusion models to revolutionize various fields, including computer vision, natural language processing, and healthcare, is vast, and further research is needed to fully realize their potential.


\section{Foundations of Alignment in Diffusion Models}



\subsection{Mathematical Foundations of Diffusion Models}


The mathematical foundations of diffusion models are rooted in the concepts of score-based models and denoising diffusion processes. These models have gained significant attention in recent years due to their ability to generate high-quality samples in various domains, including image and text generation \cite{doi:10.1109/tpami.2023.3261988}. At their core, diffusion models involve a forward process that gradually adds noise to the input data and a reverse process that learns to remove this noise to recover the original data \cite{doi:10.1145/3626235}. 

The score-based approach to diffusion models, as seen in \cite{arxiv:2208.11970}, formulates the reverse process as a sequence of noise-conditional score functions. This formulation enables the model to learn the gradient of the log-density of the data distribution, allowing for efficient sampling and generation of new data points. In contrast, denoising diffusion models, such as those described in \cite{arxiv:2303.07576}, directly model the reverse process as a sequence of denoising steps, each of which removes a small amount of noise from the input data. 

Both score-based and denoising diffusion models rely on the concept of a Markov chain, which defines a sequence of states where each state depends only on the previous state \cite{doi:10.5121/ijfcst.2015.5501}. The forward process in diffusion models can be viewed as a Markov chain that progressively adds noise to the input data, while the reverse process can be seen as a Markov chain that progressively removes this noise. This mathematical structure enables diffusion models to capture complex distributions and generate high-quality samples \cite{arxiv:2402.16369}. 

Recent advances in diffusion models have led to the development of more efficient and effective algorithms for training and sampling from these models. For example, \cite{arxiv:2210.09292} discusses various techniques for accelerating the sampling process in diffusion models, including the use of learned noise schedules and conditional sampling. Additionally, \cite{doi:10.1109/tkde.2024.3361474} provides an overview of the current state of generative diffusion models, including their applications, strengths, and limitations. 

The alignment of diffusion models with specific preferences or criteria is a critical aspect of their application in real-world domains. This alignment can be achieved through various methods, including reinforcement learning from human feedback \cite{arxiv:2305.13301} and preference-based optimization \cite{arxiv:2601.04300}. These methods enable diffusion models to generate samples that not only reflect the underlying data distribution but also meet specific requirements or preferences. 

In conclusion, the mathematical foundations of diffusion models provide a powerful framework for generative modeling and alignment with specific preferences or criteria. By understanding the underlying mathematical structure of these models, researchers and practitioners can develop more effective algorithms and applications for diffusion models, leading to significant advances in fields such as computer vision, natural language processing, and beyond \cite{doi:10.1111/cgf.15063}. As research in this area continues to evolve, we can expect to see even more innovative applications of diffusion models in the future, driven by advances in our understanding of their mathematical foundations \cite{doi:10.1109/tpami.2025.3545047}.


\subsection{Guidance Methods for Aligning Diffusion Models}


Guidance methods are a crucial component in the alignment of diffusion models with specific preferences or criteria, building on the mathematical foundations and theoretical perspectives discussed earlier. These methods enable the control of the generation process, allowing for the production of high-quality samples that meet desired attributes. In this subsection, we delve into various guidance methods, including classifier-free guidance, reinforcement learning from human feedback, and preference-based optimization, highlighting their strengths, limitations, and trade-offs, and exploring how they can be used to improve the alignment of diffusion models.

Classifier-free guidance, as introduced in \cite{arxiv:2102.09672}, has become a fundamental tool in modern diffusion models for text-guided generation. This approach offers a way to control the generation process without explicit classification, providing significant improvements in sample quality for text-to-image synthesis. By leveraging classifier-free guidance, diffusion models can be aligned with specific preferences or criteria, such as generating images that meet certain aesthetic or semantic requirements.

Reinforcement learning from human feedback (RLHF) is another guidance method that has gained significant attention. RLHF involves training a reward model on human feedback and using this model to fine-tune the diffusion model, aiming to improve alignment with human preferences \cite{arxiv:2305.13301}. This approach has shown promise in aligning diffusion models with complex and nuanced human preferences, and can be used in conjunction with other guidance methods to further improve alignment.

Preference-based optimization methods have also been explored for aligning diffusion models with user preferences. These methods incorporate user preferences into the generation process, often using preference data or reward functions. By optimizing the diffusion model to maximize user preference, these methods can improve the overall quality and relevance of the generated samples.

Theoretical insights into guidance methods have also been provided, with studies analyzing the influence of guidance on diffusion models \cite{arxiv:2403.01639}. These insights can inform the design of more effective guidance methods, and can help to identify potential limitations and challenges in the alignment of diffusion models.

Recent advances in guidance methods have focused on improving the efficiency and effectiveness of the alignment process. For instance, \cite{doi:10.1109/access.2024.3436698} proposed a novel method for fast sampling of diffusion models using Maximum Mean Discrepancy (MMD) finetuning, which significantly reduces the required number of inference timesteps. Additionally, \cite{arxiv:2404.04465} introduced a approach for aligning text-to-image diffusion models by formulating the alignment objective as the maximization of expected human utility, which unlocks the potential of leveraging readily available per-image binary signals.

As we move forward in the development of guidance methods, emerging trends include the integration of reinforcement learning and preference-based optimization, as well as the exploration of new applications such as multimodal generation and conditional sampling \cite{arxiv:2410.00083}. These advances have the potential to further improve the alignment of diffusion models with specific preferences or criteria, and can enable the application of diffusion models in a wide range of domains.

In conclusion, guidance methods play a vital role in aligning diffusion models with specific preferences or criteria, enabling the production of high-quality samples that meet desired attributes. Classifier-free guidance, reinforcement learning from human feedback, and preference-based optimization are prominent guidance methods, each with their strengths and limitations. By understanding the theoretical foundations and limitations of these methods, and by exploring new advances and applications, we can continue to improve the alignment of diffusion models, leading to more powerful and versatile models that can be applied in a wide range of domains \cite{arxiv:2408.10207}.


\subsection{Theoretical Perspectives on Diffusion Model Alignment}


Theoretical perspectives on diffusion model alignment offer a rich framework for understanding the intricacies of aligning diffusion models with specific preferences or criteria. Recent studies, such as \cite{doi:10.1145/3626235}, have highlighted the importance of diffusion models in generative tasks, including image and text synthesis. However, the alignment of these models with human preferences or specific objectives remains a challenging task. To address this challenge, researchers have explored various approaches, including reinforcement learning \cite{arxiv:2305.13301} and preference-based optimization \cite{doi:10.1007/978-3-031-72970-6_2}.

One key aspect of diffusion model alignment is the optimization landscape, which has been analyzed in \cite{arxiv:2403.01639}. This study provides valuable insights into the behavior of diffusion models under different guidance schemes, including classifier-free guidance \cite{arxiv:2207.12598}. Theoretical analysis of the optimization landscape can help identify the most effective guidance strategies and optimize the alignment process. Furthermore, the connection between diffusion models and control theory has been explored in \cite{arxiv:2301.10677}, which highlights the potential for new guidance methods inspired by control-theoretic principles.

Information-theoretic foundations of diffusion model alignment have also been investigated, with a focus on mutual information and entropy \cite{doi:10.1109/tkde.2024.3361474}. These foundations provide a theoretical basis for evaluating the alignment quality of diffusion models and can inform the design of more effective alignment methods. Additionally, the concept of Pareto alignment has been introduced in \cite{arxiv:2402.02030}, which enables the alignment of diffusion models with multiple, potentially conflicting objectives.

Recent advances in diffusion model alignment have also explored the use of multimodal evaluation metrics \cite{arxiv:2409.05033} and the development of more efficient optimization methods \cite{arxiv:2402.16359}. These advances have the potential to improve the alignment of diffusion models with human preferences and specific objectives, enabling more effective and efficient generative modeling.

Theoretical perspectives on diffusion model alignment also highlight the importance of understanding the limitations and challenges of current approaches. For example, the issue of mode collapse in diffusion models can significantly impact the alignment quality. To address this challenge, researchers have proposed various techniques, including the use of different noise schedules \cite{arxiv:2502.12154} and the development of more robust evaluation metrics.

In conclusion, theoretical perspectives on diffusion model alignment offer a valuable framework for understanding the complexities of aligning diffusion models with specific preferences or criteria. By exploring the connections between diffusion models and other areas of machine learning, such as optimization and control theory, researchers can develop more effective alignment methods and improve the overall performance of diffusion models. As the field continues to evolve, it is likely that new theoretical insights and advances in diffusion model alignment will emerge, enabling more efficient and effective generative modeling \cite{arxiv:2502.06805}. Future research directions may include the development of more robust evaluation metrics, the exploration of new guidance methods inspired by control-theoretic principles \cite{arxiv:2301.10677}, and the application of diffusion models to novel domains and tasks \cite{doi:10.1145/3626235}.


\subsection{Evaluating and Comparing Alignment Methods}


Evaluating and comparing alignment methods for diffusion models is a crucial step in understanding their strengths and weaknesses, as well as identifying areas for improvement. As noted in \cite{doi:10.1145/3626235}, the choice of alignment method can significantly impact the performance of diffusion models. Recent studies, such as \cite{doi:10.1109/tkde.2024.3361474} and \cite{arxiv:2502.06805}, have highlighted the importance of alignment in diffusion models, emphasizing the need for a comprehensive evaluation framework. This need is further underscored by the theoretical perspectives on diffusion model alignment discussed earlier, which provide a foundation for understanding the complexities of aligning diffusion models with specific preferences or criteria.

One key aspect of evaluating alignment methods is assessing their ability to optimize the diffusion process. \cite{arxiv:2305.13301} and \cite{arxiv:2402.16359} have demonstrated the effectiveness of reinforcement learning-based approaches in fine-tuning diffusion models. However, these methods can be computationally expensive and may require significant expertise in reinforcement learning. In contrast, \cite{arxiv:2510.00502} and \cite{arxiv:2507.07510} have proposed alternative approaches that leverage variational inference and divergence minimization, respectively, to align diffusion models with desired preferences. These approaches offer promising directions for improving the efficiency and effectiveness of alignment methods.

Another important consideration in evaluating alignment methods is their ability to balance competing objectives, such as sample quality and diversity. \cite{arxiv:2601.15968} and \cite{arxiv:2502.01667} have introduced hypernetwork-based and intermediate-step preference ranking approaches, respectively, to address this challenge. These methods have shown promise in improving the alignment of diffusion models with human preferences while maintaining sample diversity. Furthermore, the development of more nuanced evaluation metrics, such as those discussed in \cite{arxiv:2410.00083} and \cite{arxiv:2303.07576}, is essential for accurately assessing the performance of alignment methods.

The evaluation of alignment methods also requires careful consideration of the metrics used to assess their performance. Common metrics used in evaluating alignment methods include inception scores, Frechet inception distances, and peak signal-to-noise ratios. However, as noted in \cite{doi:10.3390/e25040633}, the choice of metric can significantly impact the evaluation results, emphasizing the need for a nuanced understanding of the strengths and limitations of each metric. Emerging trends in evaluating and comparing alignment methods include the use of multimodal evaluation metrics and the development of more efficient optimization algorithms. \cite{arxiv:2302.10907} and \cite{doi:10.1093/nsr/nwae348} have highlighted the potential of diffusion models in bioinformatics and generative AI, respectively, emphasizing the need for more efficient and effective alignment methods.

As the field of diffusion model alignment continues to evolve, it is likely that new trends and challenges will emerge, emphasizing the need for ongoing research and development in this area. The insights from existing studies, such as \cite{arxiv:2410.02543} and \cite{arxiv:2406.01572}, will play a crucial role in informing the development of more effective alignment methods. By synthesizing these insights and exploring new directions, such as \cite{arxiv:2512.15923}, researchers can unlock the full potential of diffusion models and address the open challenges and future directions in the field, including those discussed in the subsequent section. Ultimately, the refinement of evaluation and comparison methods for alignment techniques will be essential for advancing the state-of-the-art in diffusion model alignment and enabling more efficient and effective generative modeling.


\subsection{Open Challenges and Future Directions}


The alignment of diffusion models is a rapidly evolving field, with significant advancements in recent years. However, several open challenges and future directions remain to be explored. One of the primary challenges is the development of more efficient optimization methods \cite{arxiv:2502.06805}. Current optimization techniques often require substantial computational resources and time, limiting the scalability of diffusion models. To address this, researchers are exploring novel optimization algorithms and techniques, such as stochastic gradient descent and adaptive learning rates \cite{doi:10.1109/cvpr52734.2025.02198}.

Another significant challenge is the exploration of new applications for diffusion models. While they have been successfully applied to image and text generation, their potential in other domains, such as video and audio processing, remains largely unexplored \cite{doi:10.1145/3626235}. The development of diffusion models for these domains could lead to significant breakthroughs in areas like video synthesis and audio processing. Furthermore, the integration of diffusion models with other machine learning paradigms, such as reinforcement learning and transfer learning, could lead to even more powerful and flexible models \cite{arxiv:2502.10458}.

In addition to these challenges, there is a need for more comprehensive evaluation frameworks for diffusion models. Current evaluation metrics often focus on specific aspects of model performance, such as image quality or text coherence, but neglect other important factors like diversity and creativity \cite{arxiv:2508.09937}. The development of more nuanced and multi-faceted evaluation frameworks could provide a more accurate assessment of model performance and guide future research directions.

The alignment of diffusion models with human preferences and values is another critical area of research. Recent studies have demonstrated the potential of reinforcement learning and preference-based optimization for aligning diffusion models with human preferences \cite{arxiv:2404.04465}. However, these approaches often require significant amounts of human feedback and annotation, which can be time-consuming and expensive. The development of more efficient and scalable methods for aligning diffusion models with human preferences could have a significant impact on the field.

Emerging trends in diffusion model research include the development of multimodal diffusion models, which can process and generate multiple forms of data, such as images, text, and audio \cite{arxiv:2503.14504}. These models have the potential to enable more flexible and expressive generation capabilities, but also pose significant challenges in terms of training and evaluation. Another trend is the exploration of diffusion models for specific applications, such as image editing and manipulation \cite{arxiv:2601.14056}, and video synthesis \cite{doi:10.1109/cvpr52734.2025.00750}.

In conclusion, the alignment of diffusion models is a rapidly evolving field, with significant challenges and opportunities remaining to be explored. The development of more efficient optimization methods, the exploration of new applications, and the creation of more comprehensive evaluation frameworks are critical areas of research. Additionally, the alignment of diffusion models with human preferences and values, as well as the development of multimodal diffusion models, are essential for realizing the full potential of these models. As research in this field continues to advance, we can expect to see significant breakthroughs in areas like image and video synthesis, text generation, and audio processing, with potential applications in fields like computer vision, natural language processing, and robotics \cite{arxiv:2303.07576}.


\section{Methods for Aligning Diffusion Models}



\subsection{Classifier-Free Guidance Methods}


Classifier-free guidance methods have emerged as a crucial component in modern diffusion models, particularly for text-guided generation tasks. These methods enable the control of the generation process without the need for explicit classification, offering a flexible and powerful tool for aligning diffusion models with specific preferences or criteria. As highlighted in \cite{doi:10.1145/3626235}, diffusion models have shown remarkable capabilities in generating high-quality samples, and classifier-free guidance has played a significant role in this success.

One of the key approaches in classifier-free guidance is CFG++, an enhanced classifier-free guidance method that tackles the off-manifold challenges inherent in traditional classifier-free guidance \cite{arxiv:2406.08070}. CFG++ offers significant improvements in sample quality for text-to-image generation, invertibility, and reduced mode collapse. 

Model-Guidance (MG) is another novel objective for training diffusion models that addresses and removes the need for classifier-free guidance \cite{doi:10.1109/iccv51070.2023.00713}. MG results in accelerated training and exceptional quality that parallels or surpasses concurrent diffusion models with classifier-free guidance. 

Theoretical insights into diffusion guidance have also been explored, providing a deeper understanding of the influence of guidance on diffusion models \cite{arxiv:2403.01639}. This analysis has shown that incorporating diffusion guidance not only boosts classification confidence but also diminishes distribution diversity, leading to a reduction in the differential entropy of the output distribution. 

Recent studies have also focused on the development of more efficient and effective classifier-free guidance methods. For instance, \cite{arxiv:2406.05191} presents a novel technique that applies information-theoretic principles to decompose the input text prompt into its elementary components, enabling a detailed examination of how individual tokens and their interactions shape the generated image. Additionally, \cite{arxiv:2502.09935} demonstrates that only a small percentage of diffusion model parameters influence the generation of textual content within images, and proposes several applications that benefit from localizing the layers responsible for textual content generation.

The applications of classifier-free guidance methods extend beyond text-to-image synthesis. For example, \cite{arxiv:2305.13655} proposes a novel approach that leverages large language models to enhance prompt understanding in text-to-image diffusion models. 

In conclusion, classifier-free guidance methods have revolutionized the field of diffusion models, offering a powerful tool for controlling the generation process without explicit classification. The various approaches and techniques developed in this area have demonstrated remarkable capabilities in generating high-quality samples and have been applied in a wide range of domains. As research in this area continues to evolve, we can expect to see even more innovative and effective classifier-free guidance methods that will further advance the field of diffusion models and their applications. Future directions may include the development of more efficient and effective methods for optimizing classifier-free guidance, as well as the exploration of new applications and domains where these methods can be applied. As noted in \cite{arxiv:2303.07576}, the potential of diffusion models in natural language processing is vast, and classifier-free guidance methods will likely play a crucial role in unlocking this potential.


\subsection{Reinforcement Learning from Human Feedback}


Reinforcement learning from human feedback (RLHF) has emerged as a powerful approach for aligning diffusion models with human preferences and values, building on the concepts of classifier-free guidance methods discussed earlier. This methodology involves training a reward model on human feedback and utilizing this model to fine-tune the diffusion model, with the ultimate goal of enhancing alignment with human preferences \cite{arxiv:2305.13301}. The process typically entails collecting human feedback in the form of ratings or pairwise comparisons, which are then used to train a reward model that predicts the reward for a given input \cite{arxiv:2303.07576}. This reward model is subsequently employed to guide the fine-tuning of the diffusion model, encouraging it to generate samples that are more closely aligned with human preferences.

One of the primary advantages of RLHF is its ability to incorporate human values and preferences into the generation process, allowing for more nuanced and context-dependent control over the output. This is particularly significant in applications where the generation of realistic and diverse samples is crucial, such as in text-to-image synthesis or image editing \cite{arxiv:2402.16369}. Moreover, RLHF can be used to adapt diffusion models to specific tasks or domains, enabling them to generate high-quality samples that are tailored to particular requirements or constraints \cite{doi:10.1109/tpami.2025.3545047}. By leveraging RLHF, researchers can develop diffusion models that are better aligned with human preferences, which is a crucial aspect of their development and application.

Despite its potential, RLHF also poses several challenges and limitations. One of the primary concerns is the need for high-quality human feedback, which can be time-consuming and expensive to collect. Furthermore, the reward model must be carefully designed and trained to accurately predict the reward for a given input, which can be a complex and nuanced task \cite{arxiv:2402.01965}. Additionally, the fine-tuning process can be computationally expensive and may require significant resources and expertise \cite{arxiv:2502.06805}. These challenges highlight the need for ongoing research and development in RLHF, particularly in the context of preference-based optimization methods, which will be discussed in the following section.

Recent studies have explored various approaches to addressing these challenges and improving the effectiveness of RLHF. For example, \cite{arxiv:2410.02543} proposes a novel framework for RLHF that leverages evolutionary algorithms to optimize the reward model and fine-tune the diffusion model. This approach has been shown to be highly effective in adapting diffusion models to specific tasks and domains. Another study \cite{arxiv:2501.06848} presents a general framework for inference-time scaling and steering of diffusion models, which can be used to adapt RLHF to various applications and domains. These advances demonstrate the potential of RLHF to enhance the alignment of diffusion models with human preferences and values, and highlight the need for continued research and development in this area.

In conclusion, RLHF has emerged as a powerful approach for aligning diffusion models with human preferences and values, and its potential applications are vast and varied. While it poses several challenges and limitations, recent studies have explored various approaches to addressing these concerns and improving the effectiveness of RLHF. As the field continues to evolve, it is likely that RLHF will play an increasingly important role in the development of diffusion models and their applications \cite{doi:10.1016/j.media.2023.102846}. Future research directions may include exploring new approaches to collecting and incorporating human feedback, developing more efficient and effective reward models, and adapting RLHF to various applications and domains \cite{doi:10.1111/cgf.15063}. Ultimately, the potential of RLHF to enhance the alignment of diffusion models with human preferences and values makes it an exciting and promising area of research \cite{doi:10.1093/nsr/nwae348}.


\subsection{Preference-Based Optimization Methods}


Preference-based optimization methods have emerged as a crucial approach for aligning diffusion models with human preferences, enabling the incorporation of user feedback into the generation process. These methods optimize the model's parameters to better align with human preferences, often using preference data or reward functions. As highlighted in \cite{doi:10.1145/3626235}, diffusion models have shown great promise in various applications, including image synthesis, video generation, and molecule design. However, their practical applications are often hindered by the misalignment between generated samples and human preferences.

One of the key challenges in preference-based optimization is the need for effective reward functions that can capture human preferences. \cite{arxiv:2305.13301} introduces a reinforcement learning-based approach for fine-tuning diffusion models, where the reward function is designed to maximize human preferences. This approach has shown significant improvements in aligning diffusion models with human preferences, particularly in text-to-image synthesis tasks. Similarly, \cite{arxiv:2404.04465} proposes a novel approach for aligning diffusion models by formulating the alignment objective as the maximization of expected human utility.

Another important aspect of preference-based optimization is the use of preference data. \cite{doi:10.1109/cvpr52734.2025.00750} pioneers the adaptation of Direct Preference Optimization (DPO) to video diffusion models, making several key adjustments to construct a preference score that comprehensively considers both visual quality and semantic alignment between text and videos.

Theoretical insights into preference-based optimization methods are also crucial for understanding their behavior and limitations. \cite{arxiv:2408.09000} investigates the theoretical foundations of classifier-free guidance (CFG), showing that CFG interacts differently with DDPM and DDIM, and neither sampler with CFG generates the gamma-powered distribution.

Recent advances in preference-based optimization methods have also led to the development of more efficient and effective algorithms. \cite{arxiv:2409.15761} introduces a novel algorithmic framework that unifies the study of training-free guidance, providing an efficient and effective hyper-parameter searching strategy that can be readily applied to any downstream task. \cite{arxiv:2501.18865} proposes a versatile enhancement designed to boost the performance of existing guidance methods, demonstrating that REG provides a better approximation to the optimal solution than prior guidance techniques.

In conclusion, preference-based optimization methods have become a vital component of diffusion model alignment, enabling the incorporation of user preferences into the generation process. By leveraging preference data, reward functions, and theoretical insights, these methods can effectively optimize diffusion models to align with human preferences. As the field continues to evolve, future research directions may include the development of more efficient and effective algorithms, the integration of preference-based optimization with other alignment methods, and the application of these methods to various domains and tasks. With the potential to significantly improve the performance and usability of diffusion models, preference-based optimization methods are likely to play a crucial role in shaping the future of generative modeling.


\subsection{Inference-Time Alignment Control}


Inference-time alignment control for diffusion models has garnered significant attention in recent years, as it enables dynamic adjustment of the model's output to better match desired preferences or criteria without requiring additional training. Building on the concepts of preference-based optimization methods discussed in the previous section, this subsection delves into the methodologies and techniques developed for controlling the alignment of diffusion models at inference time, exploring their strengths, limitations, and potential applications. The ability to dynamically adjust the model's output is crucial for aligning diffusion models with human preferences, and inference-time alignment control offers a flexible and efficient solution for achieving this goal.

The concept of inference-time alignment control is rooted in the idea of modifying the diffusion process to incorporate user preferences or constraints, thereby generating outputs that are more aligned with the desired criteria \cite{arxiv:2410.02543}. One approach to achieving this is through the use of reinforcement learning guidance, which adapts classifier-free guidance by combining the outputs of the base and RL fine-tuned models via a geometric average \cite{arxiv:2508.21016}. This method enables dynamic control over the alignment-quality trade-off without further training, offering a flexible and efficient solution for inference-time alignment control. Furthermore, the use of reinforcement learning guidance can be seen as a natural extension of the preference-based optimization methods, where the model is optimized to maximize human preferences.

Another notable approach is the use of diffusion-sharpening, which fine-tunes diffusion models by optimizing sampling trajectories \cite{arxiv:2502.12146}. This method demonstrates superior training efficiency with faster convergence and achieves the best inference efficiency without requiring additional noise function evaluations. The concept of training-free diffusion model alignment with sampling demons has also been explored, which guides the denoising process at inference time without backpropagation through reward functions or model retraining \cite{arxiv:2410.05760}. These methods offer promising solutions for inference-time alignment control, and their integration with preference-based optimization methods can lead to even more effective and efficient alignment techniques.

The development of hypernetworks for efficient test-time alignment of diffusion models has also shown promise \cite{arxiv:2601.15968}. HyperAlign trains a hypernetwork to dynamically generate low-rank adaptation weights, allowing the denoising trajectory to be adaptively adjusted based on input latents, timesteps, and prompts for reward-conditioned alignment. This approach significantly outperforms existing fine-tuning and test-time scaling baselines in enhancing semantic consistency and visual appeal. Moreover, the concept of direct distributional optimization for provable alignment of diffusion models has been investigated \cite{arxiv:2502.02954}. This approach formulates the problem as a generic regularized loss minimization over probability distributions and directly optimizes the distribution using the Dual Averaging method.

Recent studies have also explored the use of discrete diffusion trajectory alignment via stepwise decomposition \cite{arxiv:2507.04832}. This approach decomposes the problem into a set of stepwise alignment objectives by matching the per-step posterior, enabling efficient diffusion optimization and yielding an equivalent optimal solution under additive factorization of the trajectory reward. Furthermore, the concept of beyond binary preference alignment has been investigated, which constructs a hierarchical, fine-grained evaluation criteria with domain experts and proposes a two-stage alignment framework \cite{arxiv:2601.04300}. These methods offer a range of solutions for inference-time alignment control, and their application can be seen as a crucial step towards achieving more accurate and relevant outputs in diffusion models.

The analysis of inference-time alignment control methods reveals a complex landscape of trade-offs between efficiency, flexibility, and performance. While some approaches offer superior training efficiency and inference-time control, others provide more flexible and adaptive alignment mechanisms. The choice of method ultimately depends on the specific application and desired criteria, highlighting the need for further research and development in this area. As the field continues to evolve, it is likely that new methodologies and techniques will emerge, offering improved performance, efficiency, and flexibility for inference-time alignment control in diffusion models. The potential applications of inference-time alignment control are vast, ranging from image and video generation to protein sequence design and language modeling, and its integration with other techniques, such as reinforcement learning and hypernetworks, can lead to even more powerful and flexible generative models.

In conclusion, inference-time alignment control for diffusion models is a rapidly evolving field, with various methodologies and techniques being developed to enable dynamic adjustment of the model's output to better match desired preferences or criteria. The analysis of these methods reveals a complex landscape of trade-offs between efficiency, flexibility, and performance, highlighting the need for further research and development in this area. As the field continues to evolve, it is likely that new methodologies and techniques will emerge, offering improved performance, efficiency, and flexibility for inference-time alignment control in diffusion models, and unlocking their full potential for various applications \cite{arxiv:2503.08295,arxiv:2303.07576,arxiv:2302.10907}. The future of diffusion models holds much promise, and the development of inference-time alignment control methods will play a crucial role in shaping their applications and potential impact.


\subsection{Emerging Trends and Future Directions}


The development of methods for aligning diffusion models has been a rapidly evolving field, with significant advancements in recent years. As researchers continue to push the boundaries of what is possible with diffusion models, several emerging trends and future directions have become apparent. One such trend is the integration of new techniques, such as multimodal alignment \cite{doi:10.18653/v1/d19-1136}, which enables the alignment of diffusion models across multiple modalities, including text, images, and videos. This has the potential to enable more versatile and effective generation capabilities, as demonstrated in \cite{arxiv:2211.06235} and \cite{doi:10.1145/3680528.3687621}.

Another area of focus is the application of diffusion models to novel domains, such as video generation \cite{doi:10.1109/cvpr52734.2025.00750} and 3D modeling \cite{arxiv:2601.14056}. These applications have the potential to revolutionize various fields, including computer vision, robotics, and graphics. Furthermore, researchers are exploring the use of diffusion models in combination with other generative models, such as GANs and VAEs, to enhance their performance and alignment \cite{doi:10.1145/3626235}.

The addressing of current challenges and limitations is also a key area of focus. One of the major challenges in aligning diffusion models is the issue of mode collapse, which can result in generated samples that lack diversity. To address this, researchers have proposed various techniques, such as using different noise schedules \cite{doi:10.1186/1471-2105-12-144} and modifying the model architecture \cite{doi:10.3115/1699510.1699540}. Additionally, there is a growing interest in developing more efficient and scalable alignment methods, such as \cite{doi:10.1186/s12864-020-6569-1} and \cite{arxiv:2502.06805}.

In terms of evaluation metrics, researchers are exploring new ways to assess the quality and alignment of generated samples. For example, \cite{doi:10.1109/tvcg.2020.3028975} discusses the importance of visualizing alignment results to gain a comprehensive understanding of individual findings and overall pattern structure. Moreover, \cite{arxiv:2408.13006} highlights the need for more theoretically interpretable evaluation metrics and explicit mitigation of LLM internal inconsistency.

The use of reinforcement learning and preference-based optimization is also becoming increasingly popular in aligning diffusion models \cite{doi:10.1109/cvpr52734.2025.02198} and \cite{arxiv:2402.05749}. These methods have shown significant improvements in alignment quality and efficiency, as demonstrated in \cite{arxiv:2503.12575} and \cite{doi:10.1093/nar/gkae607}.

Finally, there is a growing interest in exploring the potential of diffusion models in real-world applications, such as image editing \cite{arxiv:2307.02421} and text-to-image synthesis \cite{arxiv:2404.03653}. These applications have the potential to revolutionize various industries, including graphics, design, and advertising.

In conclusion, the field of aligning diffusion models is rapidly evolving, with significant advancements in recent years. As researchers continue to push the boundaries of what is possible with diffusion models, several emerging trends and future directions have become apparent. These include the integration of new techniques, the application of diffusion models to novel domains, and the addressing of current challenges and limitations. By exploring these areas, researchers can develop more efficient, scalable, and effective alignment methods, enabling the widespread adoption of diffusion models in various fields and applications, as discussed in \cite{arxiv:2403.04204} and \cite{arxiv:2503.14504}.


\section{Applications of Aligned Diffusion Models}



\subsection{Introduction to Applications of Aligned Diffusion Models}


The advent of aligned diffusion models has ushered in a new era of generative modeling, characterized by the ability to produce high-quality, realistic data that meets specific criteria or conditions. This subsection delves into the diverse applications of aligned diffusion models, with a particular focus on their capacity to generate data that aligns with human preferences, values, and complex criteria. As highlighted in \cite{doi:10.1145/3626235}, diffusion models have achieved state-of-the-art performance in numerous applications, including image synthesis, video generation, and molecule design.

One of the most significant applications of aligned diffusion models is in text-to-image synthesis, where the goal is to generate images that accurately reflect the input text. Recent studies, such as \cite{doi:10.1109/cvpr52729.2023.00977} and \cite{doi:10.1109/iccv51070.2023.00713}, have demonstrated the potential of aligned diffusion models in this domain. These models have shown impressive results in generating high-quality images that capture the nuances of the input text, including objects, scenes, and actions. For instance, \cite{arxiv:2305.13655} presents a novel approach that leverages large language models to enhance prompt understanding and generate more accurate images.

Another critical application of aligned diffusion models is in image editing and manipulation. As discussed in \cite{doi:10.1007/s41095-023-0371-3}, diffusion models can be used to generate high-quality images that adopt the style of a given exemplar. This has significant implications for various applications, including graphics, design, and computer vision. Furthermore, \cite{arxiv:2404.03653} proposes a novel approach that aligns text-to-image diffusion models with image-to-text concept matching, enabling more accurate and efficient image editing and manipulation.

In addition to text-to-image synthesis and image editing, aligned diffusion models have also shown promise in text generation. As highlighted in \cite{arxiv:2212.10325}, diffusion models can be used to generate high-quality text that captures the nuances of language and context. This has significant implications for various applications, including language translation, text summarization, and chatbots. Moreover, \cite{arxiv:2303.07576} provides a comprehensive overview of the applications of diffusion models in natural language processing, including text generation, language modeling, and machine translation.

The applications of aligned diffusion models are not limited to vision and language tasks. As discussed in \cite{arxiv:2410.00083}, diffusion models have also shown promise in solving inverse problems, such as image reconstruction and denoising. Furthermore, \cite{doi:10.1109/tpami.2025.3545047} provides a comprehensive overview of the applications of diffusion models in low-level vision tasks, including image denoising, deblurring, and super-resolution.

In conclusion, aligned diffusion models have shown significant promise in various applications, including text-to-image synthesis, image editing and manipulation, text generation, and inverse problems. As highlighted in \cite{arxiv:2402.16369}, the potential of diffusion models lies in their ability to generate high-quality, realistic data that meets specific criteria or conditions. However, there are still significant challenges to be addressed, including the need for more efficient and effective training methods, the development of more robust evaluation metrics, and the exploration of new applications and domains. As research in this area continues to evolve, it is likely that aligned diffusion models will play an increasingly important role in shaping the future of generative modeling and artificial intelligence. Future studies, such as \cite{arxiv:2404.07946} and \cite{arxiv:2406.05191}, will be crucial in advancing our understanding of diffusion models and their applications, and in unlocking their full potential.


\subsection{Text-to-Image Synthesis with Aligned Diffusion Models}


The application of aligned diffusion models in text-to-image synthesis has revolutionized the field of computer vision, enabling the generation of high-quality images that accurately reflect the input text. As discussed in the previous section, aligned diffusion models have shown significant promise in various applications, including text-to-image synthesis, image editing and manipulation, text generation, and inverse problems. Recent advances in diffusion models, such as those presented in \cite{doi:10.1109/tpami.2023.3261988} and \cite{arxiv:2402.16369}, have demonstrated the potential of these models in producing realistic and diverse images. The key to achieving this lies in the alignment of the diffusion model with the input text, which ensures that the generated image accurately captures the essence of the text.

One of the primary approaches to achieving alignment in text-to-image synthesis is through the use of classifier-free guidance, as discussed in \cite{arxiv:2102.09672} and \cite{arxiv:2202.09778}. This method involves training the diffusion model to predict the input text conditioned on the image, allowing for more precise control over the generation process. Another approach is to use reinforcement learning from human feedback, as presented in \cite{arxiv:2305.13301} and \cite{doi:10.1109/cvpr52734.2025.02198}, which enables the model to learn from human preferences and adapt to specific tasks. These approaches have been shown to be effective in generating high-quality images that accurately reflect the input text, and have significant implications for various applications, including computer vision, graphics, and design.

The strengths of aligned diffusion models in text-to-image synthesis lie in their ability to generate high-quality images that accurately reflect the input text. As demonstrated in \cite{arxiv:2303.07576} and \cite{doi:10.1007/978-3-031-72970-6_2}, these models can capture complex textual structures and nuances, resulting in more realistic and coherent images. Additionally, the use of diffusion models allows for efficient sampling and generation of images, making them a promising approach for practical applications. Furthermore, the application of aligned diffusion models in text-to-image synthesis has the potential to enable new applications, such as automated image generation for advertising, entertainment, and education.

However, there are also limitations and challenges associated with aligned diffusion models in text-to-image synthesis. One of the primary challenges is the need for large amounts of high-quality training data, as discussed in \cite{arxiv:2502.06805} and \cite{arxiv:2408.10207}. Additionally, the alignment process can be computationally expensive and requires careful tuning of hyperparameters, as noted in \cite{arxiv:2106.00132} and \cite{arxiv:2403.03852}. Despite these challenges, the application of aligned diffusion models in text-to-image synthesis has significant potential for future research and development.

As the field continues to evolve, it is likely that we will see the development of more efficient and effective alignment methods, as well as the integration of diffusion models with other generative models, leading to even more realistic and diverse image generation capabilities. Emerging trends, such as the use of multimodal diffusion models and the integration of diffusion models with other generative models, offer promising avenues for exploration. As noted in \cite{doi:10.1111/cgf.15063} and \cite{doi:10.1109/tpami.2025.3545047}, the development of more efficient and effective alignment methods, such as those using reinforcement learning and classifier-free guidance, is crucial for advancing the field. The following section will discuss the application of aligned diffusion models in image editing and manipulation, and how these models can be used to enable fine-grained control over image features and enable a range of tasks from simple editing operations to complex manipulations. Further studies, such as those presented in \cite{arxiv:2410.02543} and \cite{arxiv:2510.10854}, will be essential in advancing our understanding of aligned diffusion models and their applications in text-to-image synthesis.


\subsection{Image Editing and Manipulation with Aligned Diffusion Models}


The application of aligned diffusion models in image editing and manipulation has shown tremendous promise, offering fine-grained control over image features and enabling a range of tasks from simple editing operations to complex manipulations. Recent studies, such as \cite{doi:10.1109/iccv51070.2023.00685}, have demonstrated the potential of diffusion models in image editing, allowing for precise control over object placement and appearance. Similarly, \cite{arxiv:2307.02421} has introduced a novel method for drag-style manipulation, enabling users to edit images by dragging objects or regions to desired locations.

One of the key advantages of aligned diffusion models in image editing is their ability to maintain consistency and coherence in the edited images. For example, \cite{arxiv:2406.07540} has shown that diffusion models can be used to control both the structure and appearance of generated images, allowing for precise editing and manipulation. Furthermore, \cite{arxiv:2407.10592} has demonstrated the potential of diffusion models in preserving the identity and structural characteristics of objects during editing, ensuring that the edited images remain realistic and coherent.

In comparison to other approaches, such as \cite{arxiv:2207.12598}, aligned diffusion models offer improved control and flexibility in image editing. For instance, \cite{arxiv:2406.08070} has shown that manifold-constrained classifier-free guidance can improve the quality and diversity of generated images, while \cite{arxiv:2408.09000} is not supported by the provided papers and \cite{arxiv:2501.18865} has demonstrated the potential of rectified gradient guidance in enhancing the performance of diffusion models in image editing tasks.

Despite the advances in aligned diffusion models, there are still several challenges and limitations to be addressed. For example, \cite{doi:10.1109/tvcg.2024.3506992} has highlighted the issue of concept bleeding in multi-concept image generation, where the model may struggle to maintain consistency and coherence across different concepts. Additionally, \cite{arxiv:2510.17136} has demonstrated the potential of in-situ autoguidance in improving the performance of diffusion models, but also noted the need for further research in this area.

In terms of emerging trends and future directions, there is a growing interest in exploring the potential of aligned diffusion models in more complex and nuanced image editing tasks, such as \cite{doi:10.1109/tip.2026.3654412}. This approach has shown promise in offering more flexible and efficient solutions for controllable image generation, and is likely to be an active area of research in the coming years. Furthermore, the integration of aligned diffusion models with other techniques, such as \cite{arxiv:2501.18865}, is likely to lead to further improvements in image editing and manipulation capabilities.

Overall, the application of aligned diffusion models in image editing and manipulation has shown significant promise, offering fine-grained control over image features and enabling a range of tasks from simple editing operations to complex manipulations. While there are still several challenges and limitations to be addressed, the potential of aligned diffusion models in this area is substantial, and is likely to be an active area of research in the coming years. As noted in \cite{doi:10.1145/3626235}, the development of more efficient and effective diffusion models is crucial for advancing the field, and the application of aligned diffusion models in image editing and manipulation is a key area of focus. Additionally, \cite{doi:10.1109/tkde.2024.3361474} has highlighted the importance of understanding the theoretical foundations of diffusion models, and \cite{arxiv:2306.04542} has provided a comprehensive overview of the design fundamentals of diffusion models, including their application in image editing and manipulation.


\subsection{Text Generation with Aligned Diffusion Models}


The application of aligned diffusion models in text generation has garnered significant attention in recent years, owing to their potential for producing coherent and contextually accurate text \cite{arxiv:2303.07576}. Building on the success of diffusion models in image editing and manipulation, aligned diffusion models can be utilized for a range of tasks, from simple text generation to complex dialogue systems, and have shown promise in capturing long-range dependencies and contextual relationships in text data \cite{doi:10.1145/3626235}. The core idea behind aligned diffusion models for text generation is to leverage the power of diffusion processes to model the underlying structure of text sequences, and then use this structure to generate new, coherent text samples.

One of the key benefits of aligned diffusion models for text generation is their ability to capture complex patterns and relationships in text data \cite{doi:10.1109/tkde.2024.3361474}. By using a diffusion-based approach, these models can learn to represent text sequences in a more nuanced and context-dependent manner, which can lead to more realistic and engaging generated text \cite{arxiv:2302.10907}. Additionally, aligned diffusion models can be fine-tuned for specific tasks, such as text classification or sentiment analysis, by incorporating additional loss functions or constraints \cite{arxiv:2305.13301}. This flexibility and adaptability make aligned diffusion models an attractive choice for a wide range of natural language processing tasks.

Recent studies have explored the use of aligned diffusion models for text generation in various domains, including language modeling \cite{doi:10.3389/frobt.2025.1606247}, text summarization \cite{arxiv:2407.01392}, and dialogue systems \cite{arxiv:2406.01572}. These studies have demonstrated the potential of aligned diffusion models for generating high-quality, coherent text that is tailored to specific tasks or applications. For example, \cite{arxiv:2406.05191} proposed a novel approach for interpreting diffusion models using partial information decomposition, which can be used to analyze and improve the performance of aligned diffusion models for text generation.

Despite the promise of aligned diffusion models for text generation, there are still several challenges and limitations that need to be addressed \cite{arxiv:2404.07946}. One of the key challenges is the need for large amounts of training data, which can be difficult to obtain for certain tasks or domains \cite{arxiv:2502.06805}. Additionally, aligned diffusion models can be computationally expensive to train and evaluate, which can limit their applicability in certain scenarios \cite{arxiv:2305.10924}. Furthermore, the evaluation of aligned diffusion models for text generation can be challenging, as it requires the development of robust and reliable metrics for assessing the quality and coherence of generated text.

To address these challenges, researchers have proposed several techniques for improving the efficiency and effectiveness of aligned diffusion models for text generation \cite{arxiv:2303.07576}. These techniques include the use of knowledge distillation \cite{arxiv:2502.12146}, pruning \cite{arxiv:2305.10924}, and quantization \cite{arxiv:2502.06805}. Additionally, researchers have explored the use of alternative architectures and training methods, such as graph-based diffusion models \cite{arxiv:2402.16302} and reinforcement learning-based approaches \cite{arxiv:2305.13301}. As the field continues to evolve, we can expect to see further innovations and advancements in the application of aligned diffusion models for text generation, which will have significant implications for a range of natural language processing tasks and applications \cite{arxiv:2303.07576}.

In conclusion, aligned diffusion models have shown significant promise for text generation, owing to their ability to capture complex patterns and relationships in text data \cite{doi:10.1145/3626235}. While there are still several challenges and limitations that need to be addressed, researchers have proposed several techniques for improving the efficiency and effectiveness of these models \cite{arxiv:2502.06805}. The applications of aligned diffusion models in text generation will be further explored in the context of multimodal data generation and manipulation, where these models can be used to integrate and transform data across various modalities, such as text, images, and videos \cite{doi:10.1109/tkde.2024.3361474}. Future research directions may include the exploration of new architectures and training methods, the development of more robust and reliable evaluation metrics, and the application of aligned diffusion models to new and emerging domains \cite{doi:10.1109/tkde.2024.3361474}.


\subsection{Multimodal Applications of Aligned Diffusion Models}


The advent of aligned diffusion models has ushered in a new era of possibilities in the realm of multimodal data generation and manipulation. By leveraging the power of diffusion processes, these models can seamlessly integrate and transform data across various modalities, such as text, images, and videos \cite{doi:10.1145/3626235}. This subsection delves into the multifaceted applications of aligned diffusion models, exploring their capabilities in text-to-image synthesis, image-to-text generation, and other multimodal tasks.

One of the most significant advantages of aligned diffusion models is their ability to perform text-to-image synthesis \cite{doi:10.1109/cvpr52729.2023.02154}. This task involves generating images from textual descriptions, which can be achieved through the use of diffusion models that align text and image features \cite{arxiv:2404.03653}. For instance, \cite{arxiv:2211.06235} presents a coarse-to-fine alignment diffusion framework for controllable text-driven person image generation. This approach enables the generation of high-quality images that accurately reflect the input text, with applications in fields such as computer vision, graphics, and design.

In addition to text-to-image synthesis, aligned diffusion models can also be used for image-to-text generation \cite{doi:10.18653/v1/d19-1136}. This task involves generating textual descriptions from images, which can be achieved through the use of diffusion models that align image and text features \cite{doi:10.18653/v1/2021.acl-long.369}. For example, \cite{doi:10.1109/iccv51070.2023.00685} proposes a training-free method for controlling objects and contexts in synthesized images, using box constraints to guide the diffusion process.

Furthermore, aligned diffusion models can be applied to various other multimodal tasks, such as video generation and editing \cite{doi:10.1145/3680528.3687621}. This involves generating videos from textual descriptions or editing existing videos to conform to specific textual requirements \cite{arxiv:2307.02421}. For instance, \cite{doi:10.1145/3581783.3613806} presents a style-guided diffusion model for fashion synthesis, which can generate high-quality fashion images that reflect specific styles and attributes.

The applications of aligned diffusion models are not limited to the aforementioned tasks. They can also be used for data augmentation, domain adaptation, and transfer learning \cite{arxiv:2502.06805}. For example, \cite{arxiv:2404.18928} proposes an automatic adapter selection method for diffusion models, which can improve the efficiency and effectiveness of diffusion-based image synthesis. Additionally, \cite{arxiv:2402.05749} presents a generalized preference optimization framework for offline alignment, which can be used to align diffusion models with human preferences and values.

In conclusion, the multimodal applications of aligned diffusion models have the potential to revolutionize various fields, from computer vision and graphics to natural language processing and beyond. By leveraging the power of diffusion processes, these models can seamlessly integrate and transform data across various modalities, enabling a wide range of applications and use cases \cite{arxiv:2502.10458}. As research in this area continues to evolve, we can expect to see even more innovative and impactful applications of aligned diffusion models in the future \cite{arxiv:2508.06592}. With the ability to generate and manipulate multiple forms of data, aligned diffusion models are poised to play a vital role in shaping the future of artificial intelligence and multimodal data analysis \cite{arxiv:2503.14504}.


\subsection{Evaluation and Comparison of Aligned Diffusion Models}


The evaluation and comparison of aligned diffusion models are crucial steps in assessing their performance and identifying areas for improvement, building upon the multifaceted applications of these models in integrating and transforming data across various modalities, such as text, images, and videos, as discussed in the previous section. Recent advances in diffusion models have enabled the development of more efficient and effective evaluation metrics, such as those proposed in \cite{doi:10.1109/tpami.2023.3261988} and \cite{arxiv:2502.06805}. These metrics allow for a more accurate assessment of the strengths and limitations of different aligned diffusion models, which is essential for advancing the field and identifying areas for further research.

One key aspect of evaluating aligned diffusion models is the comparison of different approaches, such as the methods proposed in \cite{doi:10.18653/v1/d19-1453} and \cite{arxiv:2404.03653}. A comparative analysis of these approaches, as seen in \cite{doi:10.1109/tkde.2024.3361474}, highlights the trade-offs between different methods and identifies areas for further research, including the development of more robust and generalizable models. The evaluation of aligned diffusion models also involves assessing their strengths and limitations, such as their applications in natural language processing, including text generation and language translation, as discussed in \cite{arxiv:2303.07576}.

However, challenges such as prompt engineering, as noted in \cite{arxiv:2211.15462}, and the importance of refining network alignment, as highlighted in \cite{doi:10.1137/1.9781611976700.20}, can affect the performance of aligned diffusion models. Emerging trends and challenges in the evaluation and comparison of aligned diffusion models include the development of more efficient and effective evaluation metrics, as well as the need for more robust and generalizable models, which can be applied to various domains and tasks, such as low-level vision tasks, as discussed in \cite{doi:10.1109/tpami.2025.3545047}, and time-series forecasting tasks, as noted in \cite{arxiv:2401.03006}.

In terms of technical details, \cite{doi:10.1145/3626235} provides a comprehensive overview of the mathematical foundations of diffusion models, including the forward and reverse processes, while \cite{arxiv:2306.04542} discusses the design fundamentals of diffusion models, including the choice of noise schedule and the architecture of the diffusion model. By understanding these technical details and the strengths and limitations of aligned diffusion models, researchers and practitioners can better evaluate and compare these models, and identify areas for further research and development, ultimately advancing the field of diffusion models and their applications, as will be discussed in the following section. Future research directions may include the development of more robust and generalizable models, as well as the application of diffusion models to new domains and tasks, such as those discussed in \cite{arxiv:2410.04738} and \cite{arxiv:2405.03150}, and the development of more efficient and effective evaluation metrics, as noted in \cite{arxiv:2410.00083}.


\section{Evaluation and Comparison of Alignment Methods}



\subsection{Overview of Evaluation Metrics}


Evaluating the alignment of diffusion models is a crucial step in ensuring that these models generate high-quality, realistic data that meets specific criteria or conditions. The choice of evaluation metrics plays a significant role in assessing the performance of diffusion models, as different metrics can capture various aspects of the generated data. In this context, \cite{doi:10.1109/tpami.2023.3261988} highlights the importance of selecting appropriate metrics based on the specific application and desired outcomes of the diffusion model. For instance, metrics such as Frechet Inception Distance (FID) and Inception Score (IS) are commonly used to evaluate the quality of generated images, while metrics like BLEU score and ROUGE score are used to evaluate the quality of generated text.

One of the key challenges in evaluating diffusion models is the trade-off between quantitative and qualitative assessments. Quantitative metrics, such as FID and IS, provide a numerical score that can be used to compare the performance of different models, but they may not always capture the nuances of human perception. On the other hand, qualitative assessments, such as human evaluations, can provide a more nuanced understanding of the generated data, but they can be time-consuming and expensive to conduct. \cite{doi:10.1109/tkde.2024.3361474} discusses the importance of using a combination of quantitative and qualitative metrics to get a comprehensive understanding of model performance.

In addition to the choice of metrics, the evaluation protocol also plays a crucial role in assessing the performance of diffusion models. \cite{arxiv:2210.09292} highlights the importance of using standardized evaluation protocols to facilitate comparison and advancement in the field. Standardized protocols can help to ensure that models are evaluated fairly and consistently, which is essential for comparing the performance of different models and identifying areas for improvement.

Recent advances in diffusion models have also led to the development of new evaluation metrics and protocols. For example, \cite{arxiv:2303.07576} discusses the use of metrics such as perplexity and BLEU score to evaluate the quality of generated text. \cite{doi:10.1111/cgf.15063} highlights the importance of using metrics such as FID and IS to evaluate the quality of generated images.

Despite the progress made in evaluating diffusion models, there are still several challenges and limitations that need to be addressed. \cite{arxiv:2410.00083} highlights the importance of developing evaluation metrics and protocols that can capture the complexities of inverse problems.

In conclusion, evaluating the alignment of diffusion models is a complex task that requires careful consideration of the choice of metrics, evaluation protocol, and challenges and limitations. By selecting appropriate metrics and using standardized evaluation protocols, researchers and practitioners can get a comprehensive understanding of model performance and identify areas for improvement. As diffusion models continue to evolve and improve, it is essential to develop new evaluation metrics and protocols that can capture their complexities and nuances. Future research directions include developing more nuanced and multi-faceted evaluation frameworks, improving the efficiency and effectiveness of evaluation protocols, and exploring the use of diffusion models in new and emerging applications, as discussed in \cite{doi:10.1145/3626235} and \cite{arxiv:2402.16369}. 

The development of new evaluation metrics and protocols is an active area of research, with recent studies such as \cite{arxiv:2406.05191} and \cite{arxiv:2403.01639} proposing new approaches to evaluating diffusion models. These studies highlight the importance of continued innovation and development in the field of diffusion model evaluation, as discussed in \cite{doi:10.1088/1742-6596/2711/1/012005} and \cite{arxiv:2408.10207}. 

Overall, the evaluation of diffusion models is a critical component of their development and application, and continued research and innovation in this area are essential for advancing the field. By building on the foundations established in \cite{doi:10.1109/tpami.2025.3545047} and \cite{doi:10.7717/peerj-cs.1905}, researchers and practitioners can develop more effective and efficient evaluation metrics and protocols, ultimately leading to improved diffusion models and their applications, as highlighted in \cite{arxiv:2401.11708} and \cite{doi:10.1109/cvpr52734.2025.02198}.


\subsection{Comparison of Guidance Methods}


The development of effective guidance methods is crucial for aligning diffusion models with specific preferences or criteria, as it enables the generation of high-quality data that meets specific conditions. Among these methods, classifier-free guidance, reinforcement learning from human feedback, and preference-based optimization have garnered significant attention, each offering a unique approach to controlling the generation process. Classifier-free guidance, as seen in \cite{arxiv:2102.09672}, allows for controlling the generation process without explicit classification, offering a versatile tool for text-guided generation. However, it can suffer from mode collapse and lack of diversity, as noted in \cite{doi:10.1145/3626235}. This limitation highlights the need for careful evaluation and selection of guidance methods, as discussed in the context of evaluating the alignment of diffusion models.

Reinforcement learning from human feedback, explored in \cite{arxiv:2305.13301}, enables the alignment of models with human preferences and values. This approach involves training a reward model on human feedback and using this model to fine-tune the diffusion model. Despite its potential, it faces challenges such as requiring large amounts of high-quality human annotations and dealing with sparse rewards, as discussed in \cite{doi:10.1109/cvpr52734.2025.02198}. The use of reinforcement learning from human feedback underscores the importance of developing effective evaluation metrics and protocols, as emphasized in \cite{arxiv:2210.09292} and \cite{doi:10.1109/tkde.2024.3361474}.

Preference-based optimization methods, such as those presented in \cite{doi:10.1145/3631116}, incorporate user preferences into the generation process. These methods aim to optimize the model's parameters to better align with human preferences, often using preference data or reward functions. However, they can be computationally expensive and may require careful tuning of hyperparameters, as highlighted in \cite{doi:10.1109/tkde.2024.3361474}. The development of preference-based optimization methods highlights the need for a balanced assessment of both image quality and text alignment, as discussed in the context of evaluating aligned diffusion models.

A comparative analysis of these guidance methods reveals that each has its strengths and weaknesses. Classifier-free guidance offers flexibility and ease of implementation but may struggle with mode collapse. Reinforcement learning from human feedback provides a direct way to incorporate human preferences but is limited by the availability and quality of feedback data. Preference-based optimization methods can lead to highly aligned models but may be computationally intensive and sensitive to hyperparameters. This comparison underscores the importance of selecting appropriate guidance methods based on the specific application and requirements, as will be discussed in the context of benchmarking and standardized evaluation protocols.

Emerging trends in guidance methods include the integration of multiple approaches, such as combining classifier-free guidance with reinforcement learning, as seen in \cite{doi:10.1007/978-3-031-72970-6_2}. Additionally, there is a growing interest in developing more efficient and scalable guidance methods, such as those utilizing \cite{arxiv:2502.06805} and \cite{arxiv:2501.06848}. The development of these methods will be crucial for unlocking the full potential of diffusion models in various applications, from image and text generation to more complex tasks such as data augmentation and style transfer.

The choice of guidance method depends on the specific application and requirements. For instance, classifier-free guidance might be preferred for tasks where diversity and flexibility are crucial, while reinforcement learning from human feedback could be more suitable for applications where aligning with human values is paramount. Preference-based optimization methods might be chosen when computational resources are ample, and high alignment accuracy is necessary. This choice highlights the need for a comprehensive understanding of the strengths and limitations of different guidance methods, as well as the importance of developing new evaluation metrics and protocols that can capture the complexities of aligned diffusion models.

In conclusion, the comparison of guidance methods for aligning diffusion models highlights the trade-offs between flexibility, diversity, computational efficiency, and alignment accuracy. As research in this area continues to evolve, it is expected that hybrid approaches and novel guidance methods will emerge, offering improved performance and efficiency. The development of these methods will be crucial for advancing the field of artificial intelligence, particularly in the context of diffusion models and their applications, as will be discussed in the following sections. Future studies should focus on addressing the current limitations of guidance methods, exploring new applications, and investigating the theoretical foundations of diffusion models, as suggested in \cite{arxiv:2208.11970} and \cite{arxiv:2410.02543}.


\subsection{Benchmarking and Standardized Evaluation Protocols}


The development and evaluation of aligned diffusion models have become increasingly important in the field of artificial intelligence, with a growing need for standardized protocols to assess their performance and facilitate comparison across different models and approaches. As highlighted in \cite{doi:10.1145/3626235}, the lack of comprehensive benchmarks and evaluation frameworks hinders the advancement of diffusion models, making it challenging to determine the most effective methods for specific applications. To address this issue, researchers have proposed various benchmarks and evaluation protocols, such as the CIFAR-10 and LSUN datasets, which have become widely adopted in the community.

One of the primary challenges in evaluating aligned diffusion models is the need for a balanced assessment of both image quality and text alignment. As discussed in \cite{arxiv:2207.12598}, existing evaluation metrics, such as Frechet Inception Distance (FID) and Inception Score (IS), focus primarily on image quality, while neglecting the importance of text alignment. To overcome this limitation, researchers have introduced new metrics, such as the CLIP score, which assesses the semantic similarity between the generated image and the input text. However, as noted in \cite{doi:10.1109/cvpr52729.2023.02154}, these metrics are not without their limitations, and further research is needed to develop more comprehensive evaluation protocols.

The importance of benchmarking and standardized evaluation protocols is further emphasized in \cite{arxiv:2306.04542}, which highlights the need for a systematic and comprehensive approach to evaluating diffusion models. The authors argue that existing evaluation frameworks are often ad hoc and lack a clear understanding of the underlying design fundamentals of diffusion models. To address this issue, they propose a taxonomy of evaluation metrics and protocols, which provides a foundation for the development of more comprehensive and standardized evaluation frameworks.

In addition to the development of new evaluation metrics and protocols, researchers have also explored the use of reinforcement learning and human feedback to improve the alignment of diffusion models. As discussed in \cite{arxiv:2305.13301}, reinforcement learning provides a powerful framework for optimizing diffusion models to maximize specific objectives, such as image quality or text alignment. However, as noted in \cite{arxiv:2401.12244}, the use of reinforcement learning requires careful consideration of the reward function and the exploration-exploitation trade-off.

The use of human feedback to evaluate and improve diffusion models is also an active area of research. As highlighted in \cite{arxiv:2407.01648}, human feedback provides a valuable source of information for evaluating the quality and alignment of generated images. However, as noted in \cite{doi:10.1109/cvpr52734.2025.00750}, collecting high-quality human feedback data can be time-consuming and expensive, making it essential to develop efficient and effective methods for leveraging human feedback in the evaluation and optimization of diffusion models.

In conclusion, the development of comprehensive benchmarks and standardized evaluation protocols is crucial for advancing the field of aligned diffusion models. As researchers continue to explore new approaches and techniques, it is essential to establish a common framework for evaluating and comparing the performance of different models and methods. By leveraging existing evaluation metrics and protocols, such as those discussed in \cite{arxiv:2502.06805} and \cite{doi:10.1109/tkde.2024.3361474}, and developing new metrics and frameworks, such as those proposed in \cite{arxiv:2406.08070} and \cite{arxiv:2501.18865}, researchers can facilitate the development of more effective and efficient diffusion models, ultimately leading to significant advancements in the field of artificial intelligence. Furthermore, the integration of reinforcement learning, human feedback, and other techniques, as discussed in \cite{arxiv:2402.16359} and \cite{arxiv:2510.17136}, will play a critical role in shaping the future of diffusion models and their applications.


\subsection{Challenges and Limitations of Evaluation Metrics}


The evaluation of aligned diffusion models poses significant challenges due to the complexities inherent in assessing generative models. One of the primary concerns is the risk of overfitting to specific evaluation metrics \cite{doi:10.1162/089120103321337421}, which can lead to models that perform exceptionally well on the chosen metric but fail to generalize to real-world applications or other evaluation criteria. This issue is exacerbated by the lack of a universally accepted evaluation framework for diffusion models, making it difficult to compare different approaches directly \cite{doi:10.1109/ichi.2017.57}. Furthermore, the choice of evaluation metric itself can significantly influence the perceived performance of a model, with different metrics potentially favoring different aspects of generative quality \cite{doi:10.1093/bioinformatics/btz576}.

Building on the challenges of evaluating aligned diffusion models, it is essential to consider the limitations of current evaluation approaches. Mode collapse is another critical challenge in the evaluation of diffusion models \cite{doi:10.1145/3626235}. This phenomenon occurs when a model generates limited variations of outputs, failing to capture the full range of possibilities in the data distribution. Evaluating the extent of mode collapse is complex and requires careful consideration of both quantitative and qualitative metrics \cite{doi:10.1109/tkde.2024.3361474}. Quantitative metrics, such as the Frechet Inception Distance (FID) and Inception Score (IS), can provide insights into the diversity and realism of generated samples \cite{arxiv:2210.09292}, but they may not fully capture the nuances of mode collapse or the model's ability to generate diverse and coherent outputs.

To address these challenges, the development of more nuanced and multi-faceted evaluation frameworks is essential for accurately assessing the performance of aligned diffusion models \cite{arxiv:2303.07576}. These frameworks should incorporate a range of metrics that evaluate different aspects of generative quality, including diversity, realism, and coherence \cite{arxiv:2106.00132}. Moreover, there is a need for evaluation metrics that can assess the alignment of diffusion models with specific preferences or criteria, such as human evaluations or task-specific metrics \cite{arxiv:2305.13301}. Human evaluations, in particular, offer a direct assessment of whether the generated outputs meet the desired criteria, but they can be time-consuming and expensive to obtain \cite{arxiv:2410.02543}.

Recent advances in evaluation metrics and frameworks have shown promise in addressing some of these challenges \cite{arxiv:2402.07802}. For instance, the use of probabilistic metrics, such as those based on the Wasserstein distance or the Kullback-Leibler divergence, can provide a more comprehensive assessment of generative quality \cite{doi:10.1098/rsos.211361}. Additionally, techniques like reinforcement learning and preference-based optimization have been explored for fine-tuning diffusion models to better align with human preferences or specific objectives \cite{arxiv:2402.16359}. These developments have the potential to improve the evaluation of aligned diffusion models, enabling more effective comparison and selection of models for various applications.

Despite these advancements, significant challenges and limitations remain in the evaluation of aligned diffusion models. The lack of standardized evaluation protocols and the high computational cost of training and evaluating these models are notable barriers \cite{arxiv:2502.06805}. Furthermore, the inherent subjectivity of human evaluations and the potential for bias in preference-based optimization methods underscore the need for continued research into more objective and robust evaluation frameworks \cite{arxiv:2407.01648}. As the field continues to evolve, it is crucial to address these challenges and develop more comprehensive, robust, and efficient methods for assessing the performance of aligned diffusion models.

In conclusion, the evaluation of aligned diffusion models is a complex and multifaceted challenge that requires careful consideration of various metrics and frameworks. While significant progress has been made in addressing some of the limitations of current evaluation approaches, further research is necessary to develop more comprehensive, robust, and efficient methods for assessing the performance of these models \cite{arxiv:2302.10907}. By synthesizing insights from across the field and leveraging advances in probabilistic metrics, reinforcement learning, and preference-based optimization, it is possible to create more effective evaluation frameworks that support the continued development and application of aligned diffusion models \cite{arxiv:2410.00083}. Ultimately, the development of such frameworks will be critical for unlocking the full potential of diffusion models in a wide range of applications, from computer vision and natural language processing to biology and healthcare \cite{doi:10.3389/frobt.2025.1606247}.


\subsection{Emerging Trends and Future Directions}


The evaluation and comparison of alignment methods for diffusion models is a rapidly evolving field, with new techniques and approaches being proposed regularly. One emerging trend is the use of multimodal evaluation metrics, which can assess the alignment of diffusion models across multiple modalities, such as text, images, and videos \cite{doi:10.18653/v1/d19-1136}. This allows for a more comprehensive understanding of the model's performance and can help identify areas for improvement. Another trend is the development of more efficient and effective guidance methods, such as classifier-free guidance and reinforcement learning from human feedback \cite{doi:10.1007/s41019-024-00272-9}.

Recent studies have also highlighted the importance of benchmarking and standardized evaluation protocols for aligned diffusion models \cite{arxiv:2508.09937}. This includes the development of new metrics and evaluation frameworks that can assess the performance of diffusion models in a more nuanced and task-specific way \cite{doi:10.1186/s12864-020-6569-1}. For example, \cite{doi:10.1007/s41019-024-00272-9} proposes a novel evaluation framework for assessing the alignment of text-to-image diffusion models, while \cite{doi:10.1109/cvpr52734.2025.00750} introduces a new metric for evaluating the alignment of video diffusion models.

The use of diffusion models in natural language processing (NLP) tasks is also an emerging trend \cite{arxiv:2211.06235}. This includes the application of diffusion models to tasks such as text generation, language translation, and question answering \cite{arxiv:2502.10458}. For example, \cite{arxiv:2404.03653} proposes a novel approach for aligning text-to-image diffusion models using image-to-text concept matching, while \cite{arxiv:2404.18928} introduces a method for selecting and composing task-specific adapters for diffusion models.

Another area of research is the development of more efficient and scalable alignment methods, such as the use of parallel processing and distributed computing \cite{arxiv:2406.06382}. This can help reduce the computational cost of alignment and make it more accessible to a wider range of applications \cite{arxiv:2503.12575}. For example, \cite{arxiv:2601.14056} proposes a novel approach for positioning objects consistently and interactively in 3D scenes using diffusion models, while \cite{doi:10.1109/tip.2026.3653198} introduces a method for aligning standard and ultra-widefield fundus images using particle diffusion matching.

In terms of future directions, there are several areas that require further research and development. One area is the development of more comprehensive and standardized evaluation protocols for aligned diffusion models \cite{arxiv:2508.09937}. This includes the development of new metrics and evaluation frameworks that can assess the performance of diffusion models in a more nuanced and task-specific way. Another area is the application of diffusion models to new and emerging tasks, such as multimodal learning and human-computer interaction \cite{arxiv:2502.10458}. For example, \cite{arxiv:2511.06222} proposes a novel approach for achieving consensus in LLM alignment using self-priority optimization, while \cite{arxiv:2503.14504} provides a comprehensive survey of methods for aligning multimodal LLMs with human preference.

Overall, the evaluation and comparison of alignment methods for diffusion models is a rapidly evolving field, with new techniques and approaches being proposed regularly. As the field continues to advance, it is likely that we will see the development of more efficient and effective guidance methods, more comprehensive and standardized evaluation protocols, and the application of diffusion models to new and emerging tasks. By providing a comprehensive and in-depth analysis of the current state of the field, we can help identify areas for improvement and provide a foundation for future research and development \cite{doi:10.1145/3626235}.


\section{Challenges and Limitations of Alignment in Diffusion Models}



\subsection{Introduction to Challenges in Aligning Diffusion Models}


The alignment of diffusion models is a complex task that poses several challenges, including mode collapse, training instability, and evaluation metrics \cite{doi:10.1109/tpami.2023.3261988}. Mode collapse refers to the phenomenon where the generated samples lack diversity, resulting in a limited range of possible outcomes \cite{doi:10.1145/3626235}. This issue can be attributed to the choice of noise schedule, model architecture, and optimization algorithms used during training \cite{arxiv:2208.11970}. Training instability is another significant challenge, which can be caused by factors such as batch size, learning rate, and model complexity \cite{arxiv:2210.09292}. Furthermore, evaluating the alignment of diffusion models is a difficult task, as it requires the use of appropriate metrics that can capture the nuances of the generated samples \cite{arxiv:2410.00083}.

Recent studies have highlighted the importance of addressing these challenges in order to improve the performance of diffusion models \cite{doi:10.1109/tkde.2024.3361474}. For instance, \cite{arxiv:2303.07576} proposes the use of classifier-free guidance to improve the diversity of generated samples, while \cite{doi:10.1111/cgf.15063} suggests the use of techniques such as exponential moving average to stabilize the training process. Additionally, \cite{arxiv:2410.00083} emphasizes the need for developing more effective evaluation metrics that can accurately assess the quality of generated samples.

Despite these efforts, the alignment of diffusion models remains a challenging task, and further research is needed to develop more effective solutions \cite{doi:10.1093/nsr/nwae348}. Emerging trends, such as the use of multimodal diffusion models \cite{doi:10.1109/iccv51070.2023.00713} and the integration of reinforcement learning with diffusion models \cite{arxiv:2305.13301}, offer promising avenues for future research. Moreover, the development of more efficient and scalable algorithms for training diffusion models \cite{arxiv:2502.06805} is crucial for enabling their widespread adoption in various applications.

Theoretical insights into the behavior of diffusion models can also provide valuable guidance for addressing the challenges associated with their alignment \cite{arxiv:2403.01639}. For example, \cite{arxiv:2402.07802} presents a mathematical framework for understanding the consistency of diffusion models, while \cite{arxiv:2406.05191} proposes a technique for analyzing the information flow in diffusion models.

In conclusion, the alignment of diffusion models is a complex task that requires careful consideration of several challenges, including mode collapse, training instability, and evaluation metrics \cite{arxiv:2408.10207}. By synthesizing the existing knowledge and identifying emerging trends and challenges, researchers can work towards developing more efficient, scalable, and effective diffusion models that can be used in a wide range of applications, from computer vision to natural language processing \cite{doi:10.1109/tpami.2025.3545047}. Ultimately, the successful alignment of diffusion models has the potential to enable significant advances in fields such as generative modeling, image and video generation, and data augmentation \cite{arxiv:2402.16369}.


\subsection{Mode Collapse and Limited Diversity in Diffusion Models}


Mode collapse is a pervasive issue in diffusion models, where the generated samples lack diversity and fail to capture the complexity of the target distribution \cite{doi:10.1109/tpami.2023.3261988}. This phenomenon occurs when the model converges to a limited subset of modes in the data distribution, resulting in a lack of variability in the generated samples \cite{doi:10.1145/3626235}. As discussed in the previous section, the alignment of diffusion models is a complex task that poses several challenges, including mode collapse, training instability, and evaluation metrics \cite{arxiv:2408.10207}. The causes of mode collapse are multifaceted and can be attributed to various factors, including the choice of noise schedule, model architecture, and training objectives \cite{arxiv:2102.09672}.

One of the primary reasons for mode collapse is the use of a fixed noise schedule, which can lead to an imbalance in the exploration-exploitation trade-off \cite{arxiv:2106.00132}. When the noise schedule is too aggressive, the model may converge too quickly to a limited subset of modes, resulting in a lack of diversity in the generated samples. On the other hand, a too conservative noise schedule can lead to slow convergence and inadequate exploration of the data distribution. To mitigate this issue, researchers have proposed various techniques, such as learning the noise schedule \cite{arxiv:2202.09778} or using adaptive noise schedules \cite{arxiv:2206.05564}. These techniques can help to improve the diversity of the generated samples and address the mode collapse issue.

In addition to modifying the noise schedule, researchers have also explored the use of modified model architectures and training objectives to address mode collapse. For example, researchers have proposed using multi-scale diffusion models \cite{arxiv:2303.07576} or incorporating additional loss terms to encourage diversity in the generated samples \cite{arxiv:2310.02279}. These techniques can help to improve the diversity of the generated samples, but may also increase the computational cost and complexity of the model. Furthermore, the use of reinforcement learning and other optimization methods can also be used to fine-tune the model and encourage the generation of diverse samples that capture the complexity of the target distribution \cite{arxiv:2305.13301}.

The mitigation of mode collapse is crucial for improving the alignment of diffusion models, as it directly affects the quality and diversity of the generated samples. By addressing this issue, researchers can develop more effective diffusion models that can capture the complexity and diversity of real-world data distributions \cite{arxiv:2502.06805}. As the field continues to evolve, the development of more sophisticated techniques for mitigating mode collapse will be essential for unlocking the full potential of diffusion models. The following section will discuss the training instability and optimization challenges faced by diffusion models, which are closely related to the mode collapse issue and require careful consideration to ensure the successful alignment of diffusion models \cite{doi:10.1145/3626235}.


\subsection{Training Instability and Optimization Challenges}


Training instability and optimization challenges are significant hurdles in the development and application of diffusion models, particularly when aiming to align these models with specific preferences or criteria \cite{doi:10.1145/3626235}. The complexity of diffusion models, combined with the nature of their training objectives, can lead to unstable training processes and suboptimal solutions \cite{arxiv:2306.04542}. Several factors contribute to these challenges, including the choice of optimizer, learning rate schedule, and regularization techniques \cite{arxiv:2502.06805}.

One of the primary reasons for training instability in diffusion models is the sensitivity of the denoising process to the noise schedule \cite{arxiv:2402.07487}. The noise schedule determines how much noise is added to the input data at each step of the diffusion process, and improper scheduling can lead to poor convergence or divergence of the training process \cite{doi:10.3758/s13428-020-01366-8}. Furthermore, the use of different optimizers and learning rate schedules can significantly impact the stability and efficiency of training \cite{arxiv:2502.11420}. For instance, adaptive optimizers like AdamW have been shown to outperform traditional optimizers like SGD in certain diffusion model training scenarios \cite{doi:10.1109/tkde.2024.3361474}.

Regularization techniques also play a crucial role in mitigating training instability and improving the optimization of diffusion models \cite{arxiv:2207.12598}. Techniques such as weight decay and dropout can help prevent overfitting and promote more stable training dynamics \cite{doi:10.1109/cvpr52729.2023.02154}. Moreover, recent studies have explored the use of novel regularization methods, such as gradient clipping and learning rate warmup, to further enhance training stability \cite{arxiv:2407.01392}.

In addition to these technical considerations, the alignment of diffusion models with human preferences and values poses significant optimization challenges \cite{doi:10.1093/nsr/nwae348}. Reinforcement learning from human feedback (RLHF) has emerged as a promising approach for aligning diffusion models with human preferences \cite{arxiv:2305.13301}. However, RLHF requires careful design of the reward function and can be computationally expensive \cite{arxiv:2401.12244}.

Despite these challenges, researchers have made significant progress in developing more efficient and stable training methods for diffusion models \cite{arxiv:2502.06805}. For example, the use of mixed precision training and gradient accumulation can help reduce the computational cost and memory requirements of diffusion model training \cite{arxiv:2403.02084}. Moreover, novel architectures and training objectives, such as the use of diffusion-based generative models with implicit neural representations, have shown great promise in improving the efficiency and stability of diffusion model training \cite{arxiv:2510.17136}.

In conclusion, the training instability and optimization challenges faced by diffusion models are complex and multifaceted \cite{arxiv:2502.12154}. Addressing these challenges requires a deep understanding of the technical considerations involved in diffusion model training, as well as the development of novel methods and techniques for promoting stability and efficiency \cite{arxiv:2409.15761}. As research in this area continues to evolve, we can expect to see significant advancements in the development of more efficient, stable, and aligned diffusion models \cite{arxiv:2402.02030}. By synthesizing insights from existing studies, such as \cite{doi:10.1145/3626235}, \cite{arxiv:2306.04542}, and \cite{arxiv:2502.06805}, and exploring new approaches, such as \cite{arxiv:2501.18865} and \cite{arxiv:2506.11039}, we can unlock the full potential of diffusion models for a wide range of applications \cite{arxiv:2407.11942}. Future research directions may include the development of more sophisticated regularization techniques, the exploration of novel optimization algorithms, and the application of diffusion models to new and challenging domains \cite{arxiv:2402.16302}.


\subsection{Evaluation Metrics and Alignment Assessment}


Evaluating the alignment and quality of diffusion models is a complex task that requires careful consideration of various metrics and methods. As \cite{doi:10.1145/3626235} notes, the choice of evaluation metrics can significantly impact the assessment of model performance. Common metrics used to evaluate diffusion models include Frechet Inception Distance (FID), Inception Score (IS), and Peak Signal-to-Noise Ratio (PSNR), which provide quantitative measures of image quality and diversity. However, these metrics have limitations, such as being sensitive to specific aspects of image quality or failing to capture the nuances of human perception.

In addition to quantitative metrics, qualitative assessments, such as human evaluation and visual inspection, play a crucial role in evaluating the alignment and quality of diffusion models. As \cite{doi:10.1109/tkde.2024.3361474} highlights, human evaluation can provide valuable insights into the aesthetic quality and semantic consistency of generated images. Moreover, visual inspection can help identify mode collapse, artifacts, and other issues that may not be captured by quantitative metrics. The importance of using a combination of metrics to evaluate diffusion models is also emphasized in \cite{arxiv:2502.06805}, which notes that different metrics can capture different aspects of model performance, and a comprehensive evaluation should consider multiple metrics to get a complete picture of model quality.

The evaluation of diffusion models is closely tied to the challenges of training instability and optimization, as discussed in the previous section. The development of effective evaluation metrics and methods can help mitigate these challenges by providing a more nuanced understanding of model performance. Furthermore, \cite{arxiv:2305.13301} highlights the potential of reinforcement learning-based methods for evaluating and fine-tuning diffusion models, which can provide a more nuanced and task-specific assessment of model performance. Recent advances in evaluation metrics and methods have also focused on addressing the challenges of aligning diffusion models with human preferences and values. As \cite{arxiv:2404.04465} notes, optimizing diffusion models to maximize human utility can lead to significant improvements in model performance and alignment.

In terms of emerging trends and future directions, the development of more efficient and effective evaluation metrics and methods is crucial for facilitating the application of diffusion models in a wide range of fields. As \cite{arxiv:2410.00083} highlights, diffusion models have the potential to solve inverse problems, such as image restoration and reconstruction, and more effective evaluation metrics and methods can help unlock this potential. \cite{arxiv:2302.10907} also emphasizes the potential of diffusion models in bioinformatics, where they can be used to model complex biological systems and generate new hypotheses. The development of more effective evaluation metrics and methods can also facilitate the alignment of diffusion models with human preferences and values, as noted in \cite{doi:10.1109/cvpr52734.2025.02198}. By leveraging the strengths of different evaluation metrics and methods, researchers can develop more accurate and effective diffusion models that contribute to the advancement of various fields, including computer vision, natural language processing, and bioinformatics. Ultimately, the evaluation of diffusion models is a critical component of their development and application, and ongoing research in this area is essential for realizing the full potential of these models.


\subsection{Recent Advances and Future Directions}


Recent advances in addressing the challenges and limitations of aligning diffusion models have led to significant improvements in their performance and applicability. One of the key emerging trends is the use of multimodal diffusion models, which can align text and image features simultaneously \cite{doi:10.18653/v1/d19-1136}. This approach has shown great promise in improving the quality and diversity of generated samples, particularly in text-to-image synthesis tasks \cite{doi:10.1109/cvpr52729.2023.02154}. Furthermore, the development of novel architectures and training methods, such as the use of classifier-free guidance and adaptive projected guidance, has enabled more efficient and effective alignment of diffusion models \cite{doi:10.3115/981574.981575}.

Another significant advancement is the integration of reinforcement learning and preference-based optimization methods to improve the alignment of diffusion models \cite{doi:10.3115/990820.990835}. These approaches have been shown to be effective in incorporating user preferences into the generation process, leading to more realistic and coherent outputs \cite{doi:10.3115/1220575.1220585}. Additionally, the use of techniques such as manifold alignment and cross-modal attention has enabled the alignment of diffusion models across multiple modalities, including text, images, and videos \cite{arxiv:1705.09655}.

The development of evaluation metrics and benchmarking frameworks has also played a crucial role in advancing the field of diffusion model alignment \cite{doi:10.1109/tvcg.2020.3028975}. The creation of standardized evaluation protocols and datasets has enabled researchers to compare and contrast different alignment methods, leading to a better understanding of their strengths and limitations \cite{arxiv:2508.09937}. Moreover, the use of explainable metrics and diverse prompt templates has facilitated the assessment of alignment quality and the identification of areas for improvement \cite{arxiv:2408.13006}.

Despite these advancements, there are still several challenges and limitations associated with aligning diffusion models. One of the primary concerns is the risk of mode collapse, which can result in generated samples that lack diversity and realism. Furthermore, the alignment of diffusion models can be computationally expensive and require significant resources, particularly when dealing with large-scale datasets \cite{arxiv:2502.06805}.

To address these challenges, researchers have proposed various techniques, including the use of distillation methods, pruning, and quantization \cite{arxiv:2404.18928}. These approaches have been shown to be effective in reducing the computational requirements of diffusion models while maintaining their performance \cite{arxiv:2303.07576}. Additionally, the development of novel architectures and training methods, such as the use of transformers and attention mechanisms, has enabled more efficient and effective alignment of diffusion models \cite{arxiv:2211.06235}.

In conclusion, recent advances in addressing the challenges and limitations of aligning diffusion models have led to significant improvements in their performance and applicability. The use of multimodal diffusion models, reinforcement learning, and preference-based optimization methods has enabled more efficient and effective alignment of diffusion models. However, there are still several challenges and limitations associated with aligning diffusion models, including the risk of mode collapse and computational expenses. To address these challenges, researchers have proposed various techniques, including distillation methods, pruning, and quantization. As the field continues to evolve, it is likely that we will see the development of even more advanced techniques for aligning diffusion models, leading to further improvements in their performance and applicability \cite{doi:10.1109/cvpr52734.2025.02198}. Future research directions may include the exploration of new architectures and training methods, the development of more advanced evaluation metrics and benchmarking frameworks, and the application of diffusion models to new and emerging domains \cite{arxiv:2503.14504}.


\section{Conclusion}


The field of alignment of diffusion models has witnessed significant advancements in recent years, with various approaches being proposed to improve the performance and efficiency of these models \cite{doi:10.1109/tpami.2023.3261988}. One of the key findings of this survey is that diffusion models can be effectively aligned with human preferences and values using reinforcement learning from human feedback \cite{arxiv:2305.13301}. This approach has been shown to be particularly effective in tasks such as text-to-image synthesis, where the goal is to generate images that are not only realistic but also align with the input text \cite{arxiv:2305.13655}.

Another important aspect of diffusion models is their ability to generate high-quality samples that are diverse and realistic \cite{doi:10.1145/3626235}. This is achieved through the use of techniques such as classifier-free guidance, which allows for the generation of samples that are conditioned on a given input text \cite{arxiv:2406.08070}. However, the use of these techniques also raises concerns about the potential for mode collapse, where the generated samples lack diversity and are limited to a specific subset of the input space \cite{doi:10.1109/tvcg.2024.3506992}.

In addition to the technical advancements, there are also ethical considerations that need to be taken into account when developing and deploying diffusion models \cite{arxiv:2402.02030}. For example, the use of diffusion models for text-to-image synthesis raises concerns about the potential for generating images that are misleading or deceptive \cite{arxiv:2404.03653}. Therefore, it is essential to develop methods that can align diffusion models with human values and preferences, while also ensuring that they are fair, transparent, and accountable \cite{arxiv:2406.05882}.

Despite the significant progress made in the field of alignment of diffusion models, there are still several challenges and limitations that need to be addressed \cite{arxiv:2502.06805}. One of the key challenges is the need for more efficient and effective methods for aligning diffusion models with human preferences and values \cite{doi:10.1109/cvpr52734.2025.02198}. Another challenge is the need for more robust and generalizable evaluation metrics that can accurately assess the performance of diffusion models \cite{arxiv:2402.16369}. By addressing these challenges and limitations, we can unlock the full potential of diffusion models and enable their widespread adoption in a variety of applications \cite{doi:10.1109/tpami.2026.3658965}. 

As noted in \cite{arxiv:2406.05191}, diffusion models have the potential to learn and represent complex visual-semantic relationships. However, this also raises concerns about the potential for generating images that are misleading or deceptive \cite{doi:10.1007/s41095-023-0371-3}. Therefore, it is essential to develop methods that can align diffusion models with human values and preferences, while also ensuring that they are fair, transparent, and accountable \cite{arxiv:2404.07946}.

In the future, we can expect to see significant advancements in the field of alignment of diffusion models, with various approaches being proposed to improve the performance and efficiency of these models \cite{arxiv:2505.22165}. One potential direction for future research is the development of more efficient and effective methods for aligning diffusion models with human preferences and values \cite{arxiv:2601.04300}. Another potential direction is the development of more robust and generalizable evaluation metrics that can accurately assess the performance of diffusion models \cite{arxiv:2408.10207}. By addressing these challenges and limitations, we can unlock the full potential of diffusion models and enable their widespread adoption in a variety of applications \cite{arxiv:2502.09935}. 

Overall, the field of alignment of diffusion models is rapidly evolving, with various approaches being proposed to improve the performance and efficiency of these models \cite{arxiv:2502.08914}. While there have been significant advancements in recent years, there are still several challenges and limitations that need to be addressed \cite{arxiv:2402.07487}. Therefore, future research should focus on developing more efficient and effective methods for aligning diffusion models with human preferences and values, while also ensuring that they are fair, transparent, and accountable \cite{arxiv:2503.05149}.

\begin{thebibliography}{136}
\setlength{\itemsep}{2pt}

\bibitem{doi:10.1145/3626235}
\newblock \emph{Diffusion Models: A Comprehensive Survey of Methods and Applications}.
\newblock DOI:10.1145/3626235, 2023.
\newblock \url{https://doi.org/10.1145/3626235}

\bibitem{arxiv:1809.01036}
\newblock \emph{A Roadmap for Robust End-to-End Alignment}.
\newblock arXiv:1809.01036, 2018.
\newblock \url{https://arxiv.org/abs/1809.01036}

\bibitem{arxiv:2406.08070}
\newblock \emph{CFG++: Manifold-constrained Classifier Free Guidance for Diffusion Models}.
\newblock arXiv:2406.08070, 2024.
\newblock \url{https://arxiv.org/abs/2406.08070}

\bibitem{arxiv:2305.13301}
\newblock \emph{Training Diffusion Models with Reinforcement Learning}.
\newblock arXiv:2305.13301, 2023.
\newblock \url{https://arxiv.org/abs/2305.13301}

\bibitem{arxiv:2502.06805}
\newblock \emph{Efficient Diffusion Models: A Survey}.
\newblock arXiv:2502.06805, 2025.
\newblock \url{https://arxiv.org/abs/2502.06805}

\bibitem{arxiv:2303.07576}
\newblock \emph{Diffusion Models in NLP: A Survey}.
\newblock arXiv:2303.07576, 2023.
\newblock \url{https://arxiv.org/abs/2303.07576}

\bibitem{doi:10.1109/iccv51070.2023.00713}
\newblock \emph{Versatile Diffusion: Text, Images and Variations All in One Diffusion Model}.
\newblock DOI:10.1109/iccv51070.2023.00713, 2023.
\newblock \url{https://doi.org/10.1109/iccv51070.2023.00713}

\bibitem{arxiv:2305.13655}
\newblock \emph{LLM-grounded Diffusion: Enhancing Prompt Understanding of Text-to-Image Diffusion Models with Large Language Models}.
\newblock arXiv:2305.13655, 2023.
\newblock \url{https://arxiv.org/abs/2305.13655}

\bibitem{doi:10.1109/tpami.2025.3545047}
\newblock \emph{Diffusion Models in Low-Level Vision: A Survey}.
\newblock DOI:10.1109/tpami.2025.3545047, 2025.
\newblock \url{https://doi.org/10.1109/tpami.2025.3545047}

\bibitem{doi:10.1109/tpami.2023.3261988}
\newblock \emph{Diffusion Models in Vision: A Survey}.
\newblock DOI:10.1109/tpami.2023.3261988, 2023.
\newblock \url{https://doi.org/10.1109/tpami.2023.3261988}

\bibitem{arxiv:2208.11970}
\newblock \emph{Understanding Diffusion Models: A Unified Perspective}.
\newblock arXiv:2208.11970, 2022.
\newblock \url{https://arxiv.org/abs/2208.11970}

\bibitem{doi:10.5121/ijfcst.2015.5501}
\newblock \emph{Probabilistic Diffusion in Random Network Graphs}.
\newblock DOI:10.5121/ijfcst.2015.5501, 2015.
\newblock \url{https://doi.org/10.5121/ijfcst.2015.5501}

\bibitem{arxiv:2402.16369}
\newblock \emph{Generative AI in Vision: A Survey on Models, Metrics and Applications}.
\newblock arXiv:2402.16369, 2024.
\newblock \url{https://arxiv.org/abs/2402.16369}

\bibitem{arxiv:2210.09292}
\newblock \emph{Efficient Diffusion Models for Vision: A Survey}.
\newblock arXiv:2210.09292, 2022.
\newblock \url{https://arxiv.org/abs/2210.09292}

\bibitem{doi:10.1109/tkde.2024.3361474}
\newblock \emph{A Survey on Generative Diffusion Models}.
\newblock DOI:10.1109/tkde.2024.3361474, 2024.
\newblock \url{https://doi.org/10.1109/tkde.2024.3361474}

\bibitem{arxiv:2601.04300}
\newblock \emph{Beyond Binary Preference: Aligning Diffusion Models to Fine-grained Criteria by Decoupling Attributes}.
\newblock arXiv:2601.04300, 2026.
\newblock \url{https://arxiv.org/abs/2601.04300}

\bibitem{doi:10.1111/cgf.15063}
\newblock \emph{State of the Art on Diffusion Models for Visual Computing}.
\newblock DOI:10.1111/cgf.15063, 2024.
\newblock \url{https://doi.org/10.1111/cgf.15063}

\bibitem{arxiv:2102.09672}
\newblock \emph{Improved Denoising Diffusion Probabilistic Models}.
\newblock arXiv:2102.09672, 2021.
\newblock \url{https://arxiv.org/abs/2102.09672}

\bibitem{arxiv:2403.01639}
\newblock \emph{Theoretical Insights for Diffusion Guidance: A Case Study for Gaussian Mixture Models}.
\newblock arXiv:2403.01639, 2024.
\newblock \url{https://arxiv.org/abs/2403.01639}

\bibitem{doi:10.1109/access.2024.3436698}
\newblock \emph{Fast Inference in Denoising Diffusion Models via MMD Finetuning}.
\newblock DOI:10.1109/access.2024.3436698, 2024.
\newblock \url{https://doi.org/10.1109/access.2024.3436698}

\bibitem{arxiv:2404.04465}
\newblock \emph{Aligning Diffusion Models by Optimizing Human Utility}.
\newblock arXiv:2404.04465, 2024.
\newblock \url{https://arxiv.org/abs/2404.04465}

\bibitem{arxiv:2410.00083}
\newblock \emph{A Survey on Diffusion Models for Inverse Problems}.
\newblock arXiv:2410.00083, 2024.
\newblock \url{https://arxiv.org/abs/2410.00083}

\bibitem{arxiv:2408.10207}
\newblock \emph{A Comprehensive Survey on Diffusion Models and Their Applications}.
\newblock arXiv:2408.10207, 2024.
\newblock \url{https://arxiv.org/abs/2408.10207}

\bibitem{doi:10.1007/978-3-031-72970-6_2}
\newblock \emph{Text-Anchored Score Composition: Tackling Condition Misalignment in Text-to-Image Diffusion Models}.
\newblock DOI:10.1007/978-3-031-72970-6\_2, 2024.
\newblock \url{https://doi.org/10.1007/978-3-031-72970-6_2}

\bibitem{arxiv:2207.12598}
\newblock \emph{Classifier-Free Diffusion Guidance}.
\newblock arXiv:2207.12598, 2022.
\newblock \url{https://arxiv.org/abs/2207.12598}

\bibitem{arxiv:2301.10677}
\newblock \emph{Imitating Human Behaviour with Diffusion Models}.
\newblock arXiv:2301.10677, 2023.
\newblock \url{https://arxiv.org/abs/2301.10677}

\bibitem{arxiv:2402.02030}
\newblock \emph{Panacea: Pareto Alignment via Preference Adaptation for LLMs}.
\newblock arXiv:2402.02030, 2024.
\newblock \url{https://arxiv.org/abs/2402.02030}

\bibitem{arxiv:2409.05033}
\newblock \emph{A Survey on Diffusion Models for Recommender Systems}.
\newblock arXiv:2409.05033, 2024.
\newblock \url{https://arxiv.org/abs/2409.05033}

\bibitem{arxiv:2402.16359}
\newblock \emph{Feedback Efficient Online Fine-Tuning of Diffusion Models}.
\newblock arXiv:2402.16359, 2024.
\newblock \url{https://arxiv.org/abs/2402.16359}

\bibitem{arxiv:2502.12154}
\newblock \emph{Diffusion Models without Classifier-free Guidance}.
\newblock arXiv:2502.12154, 2025.
\newblock \url{https://arxiv.org/abs/2502.12154}

\bibitem{arxiv:2510.00502}
\newblock \emph{Diffusion Alignment as Variational Expectation-Maximization}.
\newblock arXiv:2510.00502, 2025.
\newblock \url{https://arxiv.org/abs/2510.00502}

\bibitem{arxiv:2507.07510}
\newblock \emph{Divergence Minimization Preference Optimization for Diffusion Model Alignment}.
\newblock arXiv:2507.07510, 2025.
\newblock \url{https://arxiv.org/abs/2507.07510}

\bibitem{arxiv:2601.15968}
\newblock \emph{HyperAlign: Hypernetwork for Efficient Test-Time Alignment of Diffusion Models}.
\newblock arXiv:2601.15968, 2026.
\newblock \url{https://arxiv.org/abs/2601.15968}

\bibitem{arxiv:2502.01667}
\newblock \emph{Refining Alignment Framework for Diffusion Models with Intermediate-Step Preference Ranking}.
\newblock arXiv:2502.01667, 2025.
\newblock \url{https://arxiv.org/abs/2502.01667}

\bibitem{doi:10.3390/e25040633}
\newblock \emph{How Much Is Enough? A Study on Diffusion Times in Score-Based Generative Models}.
\newblock DOI:10.3390/e25040633, 2023.
\newblock \url{https://doi.org/10.3390/e25040633}

\bibitem{arxiv:2302.10907}
\newblock \emph{Diffusion Models in Bioinformatics: A New Wave of Deep Learning Revolution in Action}.
\newblock arXiv:2302.10907, 2023.
\newblock \url{https://arxiv.org/abs/2302.10907}

\bibitem{doi:10.1093/nsr/nwae348}
\newblock \emph{Opportunities and challenges of diffusion models for generative AI}.
\newblock DOI:10.1093/nsr/nwae348, 2024.
\newblock \url{https://doi.org/10.1093/nsr/nwae348}

\bibitem{arxiv:2410.02543}
\newblock \emph{Diffusion Models are Evolutionary Algorithms}.
\newblock arXiv:2410.02543, 2024.
\newblock \url{https://arxiv.org/abs/2410.02543}

\bibitem{arxiv:2406.01572}
\newblock \emph{Unlocking Guidance for Discrete State-Space Diffusion and Flow Models}.
\newblock arXiv:2406.01572, 2024.
\newblock \url{https://arxiv.org/abs/2406.01572}

\bibitem{arxiv:2512.15923}
\newblock \emph{A Unification of Discrete, Gaussian, and Simplicial Diffusion}.
\newblock arXiv:2512.15923, 2025.
\newblock \url{https://arxiv.org/abs/2512.15923}

\bibitem{doi:10.1109/cvpr52734.2025.02198}
\newblock \emph{Towards Better Alignment: Training Diffusion Models with Reinforcement Learning Against Sparse Rewards}.
\newblock DOI:10.1109/cvpr52734.2025.02198, 2025.
\newblock \url{https://doi.org/10.1109/cvpr52734.2025.02198}

\bibitem{arxiv:2502.10458}
\newblock \emph{I Think, Therefore I Diffuse: Enabling Multimodal In-Context Reasoning in Diffusion Models}.
\newblock arXiv:2502.10458, 2025.
\newblock \url{https://arxiv.org/abs/2502.10458}

\bibitem{arxiv:2508.09937}
\newblock \emph{A Comprehensive Evaluation framework of Alignment Techniques for LLMs}.
\newblock arXiv:2508.09937, 2025.
\newblock \url{https://arxiv.org/abs/2508.09937}

\bibitem{arxiv:2503.14504}
\newblock \emph{Aligning Multimodal LLM with Human Preference: A Survey}.
\newblock arXiv:2503.14504, 2025.
\newblock \url{https://arxiv.org/abs/2503.14504}

\bibitem{arxiv:2601.14056}
\newblock \emph{POCI-Diff: Position Objects Consistently and Interactively with 3D-Layout Guided Diffusion}.
\newblock arXiv:2601.14056, 2026.
\newblock \url{https://arxiv.org/abs/2601.14056}

\bibitem{doi:10.1109/cvpr52734.2025.00750}
\newblock \emph{VideoDPO: Omni-Preference Alignment for Video Diffusion Generation}.
\newblock DOI:10.1109/cvpr52734.2025.00750, 2025.
\newblock \url{https://doi.org/10.1109/cvpr52734.2025.00750}

\bibitem{arxiv:2406.05191}
\newblock \emph{DiffusionPID: Interpreting Diffusion via Partial Information Decomposition}.
\newblock arXiv:2406.05191, 2024.
\newblock \url{https://arxiv.org/abs/2406.05191}

\bibitem{arxiv:2502.09935}
\newblock \emph{Precise Parameter Localization for Textual Generation in Diffusion Models}.
\newblock arXiv:2502.09935, 2025.
\newblock \url{https://arxiv.org/abs/2502.09935}

\bibitem{arxiv:2402.01965}
\newblock \emph{Analyzing Neural Network-Based Generative Diffusion Models through Convex Optimization}.
\newblock arXiv:2402.01965, 2024.
\newblock \url{https://arxiv.org/abs/2402.01965}

\bibitem{arxiv:2501.06848}
\newblock \emph{A General Framework for Inference-time Scaling and Steering of Diffusion Models}.
\newblock arXiv:2501.06848, 2025.
\newblock \url{https://arxiv.org/abs/2501.06848}

\bibitem{doi:10.1016/j.media.2023.102846}
\newblock \emph{Diffusion models in medical imaging: A comprehensive survey}.
\newblock DOI:10.1016/j.media.2023.102846, 2023.
\newblock \url{https://doi.org/10.1016/j.media.2023.102846}

\bibitem{arxiv:2408.09000}
\newblock \emph{Classifier-Free Guidance is a Predictor-Corrector}.
\newblock arXiv:2408.09000, 2024.
\newblock \url{https://arxiv.org/abs/2408.09000}

\bibitem{arxiv:2409.15761}
\newblock \emph{TFG: Unified Training-Free Guidance for Diffusion Models}.
\newblock arXiv:2409.15761, 2024.
\newblock \url{https://arxiv.org/abs/2409.15761}

\bibitem{arxiv:2501.18865}
\newblock \emph{REG: Rectified Gradient Guidance for Conditional Diffusion Models}.
\newblock arXiv:2501.18865, 2025.
\newblock \url{https://arxiv.org/abs/2501.18865}

\bibitem{arxiv:2508.21016}
\newblock \emph{Inference-Time Alignment Control for Diffusion Models with Reinforcement Learning Guidance}.
\newblock arXiv:2508.21016, 2025.
\newblock \url{https://arxiv.org/abs/2508.21016}

\bibitem{arxiv:2502.12146}
\newblock \emph{Diffusion-Sharpening: Fine-tuning Diffusion Models with Denoising Trajectory Sharpening}.
\newblock arXiv:2502.12146, 2025.
\newblock \url{https://arxiv.org/abs/2502.12146}

\bibitem{arxiv:2410.05760}
\newblock \emph{Training-free Diffusion Model Alignment with Sampling Demons}.
\newblock arXiv:2410.05760, 2024.
\newblock \url{https://arxiv.org/abs/2410.05760}

\bibitem{arxiv:2502.02954}
\newblock \emph{Direct Distributional Optimization for Provable Alignment of Diffusion Models}.
\newblock arXiv:2502.02954, 2025.
\newblock \url{https://arxiv.org/abs/2502.02954}

\bibitem{arxiv:2507.04832}
\newblock \emph{Discrete Diffusion Trajectory Alignment via Stepwise Decomposition}.
\newblock arXiv:2507.04832, 2025.
\newblock \url{https://arxiv.org/abs/2507.04832}

\bibitem{arxiv:2503.08295}
\newblock \emph{D3PO: Preference-Based Alignment of Discrete Diffusion Models}.
\newblock arXiv:2503.08295, 2025.
\newblock \url{https://arxiv.org/abs/2503.08295}

\bibitem{doi:10.18653/v1/d19-1136}
\newblock \emph{Vecalign: Improved Sentence Alignment in Linear Time and Space}.
\newblock DOI:10.18653/v1/d19-1136, 2019.
\newblock \url{https://doi.org/10.18653/v1/d19-1136}

\bibitem{arxiv:2211.06235}
\newblock \emph{HumanDiffusion: a Coarse-to-Fine Alignment Diffusion Framework for Controllable Text-Driven Person Image Generation}.
\newblock arXiv:2211.06235, 2022.
\newblock \url{https://arxiv.org/abs/2211.06235}

\bibitem{doi:10.1145/3680528.3687621}
\newblock \emph{Text-Guided Texturing by Synchronized Multi-View Diffusion}.
\newblock DOI:10.1145/3680528.3687621, 2024.
\newblock \url{https://doi.org/10.1145/3680528.3687621}

\bibitem{doi:10.1186/1471-2105-12-144}
\newblock \emph{Meta-Alignment with Crumble and Prune: Partitioning very large alignment problems for performance and parallelization}.
\newblock DOI:10.1186/1471-2105-12-144, 2011.
\newblock \url{https://doi.org/10.1186/1471-2105-12-144}

\bibitem{doi:10.3115/1699510.1699540}
\newblock \emph{Improved Word Alignment with Statistics and Linguistic Heuristics}.
\newblock DOI:10.3115/1699510.1699540, 2009.
\newblock \url{https://doi.org/10.3115/1699510.1699540}

\bibitem{doi:10.1186/s12864-020-6569-1}
\newblock \emph{GSAlign: an efficient sequence alignment tool for intra-species genomes}.
\newblock DOI:10.1186/s12864-020-6569-1, 2020.
\newblock \url{https://doi.org/10.1186/s12864-020-6569-1}

\bibitem{doi:10.1109/tvcg.2020.3028975}
\newblock \emph{A Survey of Text Alignment Visualization}.
\newblock DOI:10.1109/tvcg.2020.3028975, 2020.
\newblock \url{https://doi.org/10.1109/tvcg.2020.3028975}

\bibitem{arxiv:2408.13006}
\newblock \emph{Systematic Evaluation of LLM-as-a-Judge in LLM Alignment Tasks: Explainable Metrics and Diverse Prompt Templates}.
\newblock arXiv:2408.13006, 2024.
\newblock \url{https://arxiv.org/abs/2408.13006}

\bibitem{arxiv:2402.05749}
\newblock \emph{Generalized Preference Optimization: A Unified Approach to Offline Alignment}.
\newblock arXiv:2402.05749, 2024.
\newblock \url{https://arxiv.org/abs/2402.05749}

\bibitem{arxiv:2503.12575}
\newblock \emph{BalancedDPO: Adaptive Multi-Metric Alignment}.
\newblock arXiv:2503.12575, 2025.
\newblock \url{https://arxiv.org/abs/2503.12575}

\bibitem{doi:10.1093/nar/gkae607}
\newblock \emph{SigAlign: an alignment algorithm guided by explicit similarity criteria}.
\newblock DOI:10.1093/nar/gkae607, 2024.
\newblock \url{https://doi.org/10.1093/nar/gkae607}

\bibitem{arxiv:2307.02421}
\newblock \emph{DragonDiffusion: Enabling Drag-style Manipulation on Diffusion Models}.
\newblock arXiv:2307.02421, 2023.
\newblock \url{https://arxiv.org/abs/2307.02421}

\bibitem{arxiv:2404.03653}
\newblock \emph{CoMat: Aligning Text-to-Image Diffusion Model with Image-to-Text Concept Matching}.
\newblock arXiv:2404.03653, 2024.
\newblock \url{https://arxiv.org/abs/2404.03653}

\bibitem{arxiv:2403.04204}
\newblock \emph{On the Essence and Prospect: An Investigation of Alignment Approaches for Big Models}.
\newblock arXiv:2403.04204, 2024.
\newblock \url{https://arxiv.org/abs/2403.04204}

\bibitem{doi:10.1109/cvpr52729.2023.00977}
\newblock \emph{ERNIE-ViLG 2.0: Improving Text-to-Image Diffusion Model with Knowledge-Enhanced Mixture-of-Denoising-Experts}.
\newblock DOI:10.1109/cvpr52729.2023.00977, 2023.
\newblock \url{https://doi.org/10.1109/cvpr52729.2023.00977}

\bibitem{doi:10.1007/s41095-023-0371-3}
\newblock \emph{Taming diffusion model for exemplar-based image translation}.
\newblock DOI:10.1007/s41095-023-0371-3, 2024.
\newblock \url{https://doi.org/10.1007/s41095-023-0371-3}

\bibitem{arxiv:2212.10325}
\newblock \emph{SeqDiffuSeq: Text Diffusion with Encoder-Decoder Transformers}.
\newblock arXiv:2212.10325, 2022.
\newblock \url{https://arxiv.org/abs/2212.10325}

\bibitem{arxiv:2404.07946}
\newblock \emph{Towards Faster Training of Diffusion Models: An Inspiration of A Consistency Phenomenon}.
\newblock arXiv:2404.07946, 2024.
\newblock \url{https://arxiv.org/abs/2404.07946}

\bibitem{arxiv:2202.09778}
\newblock \emph{Pseudo Numerical Methods for Diffusion Models on Manifolds}.
\newblock arXiv:2202.09778, 2022.
\newblock \url{https://arxiv.org/abs/2202.09778}

\bibitem{arxiv:2106.00132}
\newblock \emph{On Fast Sampling of Diffusion Probabilistic Models}.
\newblock arXiv:2106.00132, 2021.
\newblock \url{https://arxiv.org/abs/2106.00132}

\bibitem{arxiv:2403.03852}
\newblock \emph{Accelerating Convergence of Score-Based Diffusion Models, Provably}.
\newblock arXiv:2403.03852, 2024.
\newblock \url{https://arxiv.org/abs/2403.03852}

\bibitem{arxiv:2510.10854}
\newblock \emph{Discrete State Diffusion Models: A Sample Complexity Perspective}.
\newblock arXiv:2510.10854, 2025.
\newblock \url{https://arxiv.org/abs/2510.10854}

\bibitem{doi:10.1109/iccv51070.2023.00685}
\newblock \emph{BoxDiff: Text-to-Image Synthesis with Training-Free Box-Constrained Diffusion}.
\newblock DOI:10.1109/iccv51070.2023.00685, 2023.
\newblock \url{https://doi.org/10.1109/iccv51070.2023.00685}

\bibitem{arxiv:2406.07540}
\newblock \emph{Ctrl-X: Controlling Structure and Appearance for Text-To-Image Generation Without Guidance}.
\newblock arXiv:2406.07540, 2024.
\newblock \url{https://arxiv.org/abs/2406.07540}

\bibitem{arxiv:2407.10592}
\newblock \emph{InsertDiffusion: Identity Preserving Visualization of Objects through a Training-Free Diffusion Architecture}.
\newblock arXiv:2407.10592, 2024.
\newblock \url{https://arxiv.org/abs/2407.10592}

\bibitem{doi:10.1109/tvcg.2024.3506992}
\newblock \emph{Isolated Diffusion: Optimizing Multi-Concept Text-to-Image Generation Training-Freely With Isolated Diffusion Guidance}.
\newblock DOI:10.1109/tvcg.2024.3506992, 2024.
\newblock \url{https://doi.org/10.1109/tvcg.2024.3506992}

\bibitem{arxiv:2510.17136}
\newblock \emph{In-situ Autoguidance: Eliciting Self-Correction in Diffusion Models}.
\newblock arXiv:2510.17136, 2025.
\newblock \url{https://arxiv.org/abs/2510.17136}

\bibitem{doi:10.1109/tip.2026.3654412}
\newblock \emph{LaCon: Late-Constraint Controllable Visual Generation}.
\newblock DOI:10.1109/tip.2026.3654412, 2026.
\newblock \url{https://doi.org/10.1109/tip.2026.3654412}

\bibitem{arxiv:2306.04542}
\newblock \emph{On the Design Fundamentals of Diffusion Models: A Survey}.
\newblock arXiv:2306.04542, 2023.
\newblock \url{https://arxiv.org/abs/2306.04542}

\bibitem{doi:10.3389/frobt.2025.1606247}
\newblock \emph{Diffusion models for robotic manipulation: a survey}.
\newblock DOI:10.3389/frobt.2025.1606247, 2025.
\newblock \url{https://doi.org/10.3389/frobt.2025.1606247}

\bibitem{arxiv:2407.01392}
\newblock \emph{Diffusion Forcing: Next-token Prediction Meets Full-Sequence Diffusion}.
\newblock arXiv:2407.01392, 2024.
\newblock \url{https://arxiv.org/abs/2407.01392}

\bibitem{arxiv:2305.10924}
\newblock \emph{Structural Pruning for Diffusion Models}.
\newblock arXiv:2305.10924, 2023.
\newblock \url{https://arxiv.org/abs/2305.10924}

\bibitem{arxiv:2402.16302}
\newblock \emph{Graph Diffusion Policy Optimization}.
\newblock arXiv:2402.16302, 2024.
\newblock \url{https://arxiv.org/abs/2402.16302}

\bibitem{doi:10.1109/cvpr52729.2023.02154}
\newblock \emph{LayoutDiffusion: Controllable Diffusion Model for Layout-to-Image Generation}.
\newblock DOI:10.1109/cvpr52729.2023.02154, 2023.
\newblock \url{https://doi.org/10.1109/cvpr52729.2023.02154}

\bibitem{doi:10.18653/v1/2021.acl-long.369}
\newblock \emph{Mask-Align: Self-Supervised Neural Word Alignment}.
\newblock DOI:10.18653/v1/2021.acl-long.369, 2021.
\newblock \url{https://doi.org/10.18653/v1/2021.acl-long.369}

\bibitem{doi:10.1145/3581783.3613806}
\newblock \emph{SGDiff: A Style Guided Diffusion Model for Fashion Synthesis}.
\newblock DOI:10.1145/3581783.3613806, 2023.
\newblock \url{https://doi.org/10.1145/3581783.3613806}

\bibitem{arxiv:2404.18928}
\newblock \emph{Stylus: Automatic Adapter Selection for Diffusion Models}.
\newblock arXiv:2404.18928, 2024.
\newblock \url{https://arxiv.org/abs/2404.18928}

\bibitem{arxiv:2508.06592}
\newblock \emph{Towards Integrated Alignment}.
\newblock arXiv:2508.06592, 2025.
\newblock \url{https://arxiv.org/abs/2508.06592}

\bibitem{doi:10.18653/v1/d19-1453}
\newblock \emph{Jointly Learning to Align and Translate with Transformer Models}.
\newblock DOI:10.18653/v1/d19-1453, 2019.
\newblock \url{https://doi.org/10.18653/v1/d19-1453}

\bibitem{arxiv:2211.15462}
\newblock \emph{Investigating Prompt Engineering in Diffusion Models}.
\newblock arXiv:2211.15462, 2022.
\newblock \url{https://arxiv.org/abs/2211.15462}

\bibitem{doi:10.1137/1.9781611976700.20}
\newblock \emph{Refining Network Alignment to Improve Matched Neighborhood Consistency}.
\newblock DOI:10.1137/1.9781611976700.20, 2021.
\newblock \url{https://doi.org/10.1137/1.9781611976700.20}

\bibitem{arxiv:2401.03006}
\newblock \emph{The Rise of Diffusion Models in Time-Series Forecasting}.
\newblock arXiv:2401.03006, 2024.
\newblock \url{https://arxiv.org/abs/2401.03006}

\bibitem{arxiv:2410.04738}
\newblock \emph{Diffusion Models in 3D Vision: A Survey}.
\newblock arXiv:2410.04738, 2024.
\newblock \url{https://arxiv.org/abs/2410.04738}

\bibitem{arxiv:2405.03150}
\newblock \emph{Video Diffusion Models: A Survey}.
\newblock arXiv:2405.03150, 2024.
\newblock \url{https://arxiv.org/abs/2405.03150}

\bibitem{doi:10.1088/1742-6596/2711/1/012005}
\newblock \emph{Review: Recent advances for the diffusion model}.
\newblock DOI:10.1088/1742-6596/2711/1/012005, 2024.
\newblock \url{https://doi.org/10.1088/1742-6596/2711/1/012005}

\bibitem{doi:10.7717/peerj-cs.1905}
\newblock \emph{Diffusion models in text generation: a survey}.
\newblock DOI:10.7717/peerj-cs.1905, 2024.
\newblock \url{https://doi.org/10.7717/peerj-cs.1905}

\bibitem{arxiv:2401.11708}
\newblock \emph{Mastering Text-to-Image Diffusion: Recaptioning, Planning, and Generating with Multimodal LLMs}.
\newblock arXiv:2401.11708, 2024.
\newblock \url{https://arxiv.org/abs/2401.11708}

\bibitem{doi:10.1145/3631116}
\newblock \emph{DiffuRec: A Diffusion Model for Sequential Recommendation}.
\newblock DOI:10.1145/3631116, 2023.
\newblock \url{https://doi.org/10.1145/3631116}

\bibitem{arxiv:2401.12244}
\newblock \emph{Large-scale Reinforcement Learning for Diffusion Models}.
\newblock arXiv:2401.12244, 2024.
\newblock \url{https://arxiv.org/abs/2401.12244}

\bibitem{arxiv:2407.01648}
\newblock \emph{Aligning Target-Aware Molecule Diffusion Models with Exact Energy Optimization}.
\newblock arXiv:2407.01648, 2024.
\newblock \url{https://arxiv.org/abs/2407.01648}

\bibitem{doi:10.1162/089120103321337421}
\newblock \emph{A Systematic Comparison of Various Statistical Alignment Models}.
\newblock DOI:10.1162/089120103321337421, 2003.
\newblock \url{https://doi.org/10.1162/089120103321337421}

\bibitem{doi:10.1109/ichi.2017.57}
\newblock \emph{Evaluation of Trace Alignment Quality and its Application in Medical Process Mining}.
\newblock DOI:10.1109/ichi.2017.57, 2017.
\newblock \url{https://doi.org/10.1109/ichi.2017.57}

\bibitem{doi:10.1093/bioinformatics/btz576}
\newblock \emph{How sequence alignment scores correspond to probability models}.
\newblock DOI:10.1093/bioinformatics/btz576, 2019.
\newblock \url{https://doi.org/10.1093/bioinformatics/btz576}

\bibitem{arxiv:2402.07802}
\newblock \emph{Towards a mathematical theory for consistency training in diffusion models}.
\newblock arXiv:2402.07802, 2024.
\newblock \url{https://arxiv.org/abs/2402.07802}

\bibitem{doi:10.1098/rsos.211361}
\newblock \emph{Model comparison via simplicial complexes and persistent homology}.
\newblock DOI:10.1098/rsos.211361, 2022.
\newblock \url{https://doi.org/10.1098/rsos.211361}

\bibitem{doi:10.1007/s41019-024-00272-9}
\newblock \emph{DTIA: Disruptive Text-Image Alignment for Countering Text-to-Image Diffusion Model Personalization}.
\newblock DOI:10.1007/s41019-024-00272-9, 2024.
\newblock \url{https://doi.org/10.1007/s41019-024-00272-9}

\bibitem{arxiv:2406.06382}
\newblock \emph{Diffusion-RPO: Aligning Diffusion Models through Relative Preference Optimization}.
\newblock arXiv:2406.06382, 2024.
\newblock \url{https://arxiv.org/abs/2406.06382}

\bibitem{doi:10.1109/tip.2026.3653198}
\newblock \emph{Particle Diffusion Matching: Random Walk Correspondence Search for the Alignment of Standard and Ultra-Widefield Fundus Images}.
\newblock DOI:10.1109/tip.2026.3653198, 2026.
\newblock \url{https://doi.org/10.1109/tip.2026.3653198}

\bibitem{arxiv:2511.06222}
\newblock \emph{SPA: Achieving Consensus in LLM Alignment via Self-Priority Optimization}.
\newblock arXiv:2511.06222, 2025.
\newblock \url{https://arxiv.org/abs/2511.06222}

\bibitem{arxiv:2206.05564}
\newblock \emph{gDDIM: Generalized denoising diffusion implicit models}.
\newblock arXiv:2206.05564, 2022.
\newblock \url{https://arxiv.org/abs/2206.05564}

\bibitem{arxiv:2310.02279}
\newblock \emph{Consistency Trajectory Models: Learning Probability Flow ODE Trajectory of Diffusion}.
\newblock arXiv:2310.02279, 2023.
\newblock \url{https://arxiv.org/abs/2310.02279}

\bibitem{arxiv:2402.07487}
\newblock \emph{Score-based Diffusion Models via Stochastic Differential Equations - a Technical Tutorial}.
\newblock arXiv:2402.07487, 2024.
\newblock \url{https://arxiv.org/abs/2402.07487}

\bibitem{doi:10.3758/s13428-020-01366-8}
\newblock \emph{Improving parameter recovery for conflict drift-diffusion models}.
\newblock DOI:10.3758/s13428-020-01366-8, 2020.
\newblock \url{https://doi.org/10.3758/s13428-020-01366-8}

\bibitem{arxiv:2502.11420}
\newblock \emph{Training-Free Guidance Beyond Differentiability: Scalable Path Steering with Tree Search in Diffusion and Flow Models}.
\newblock arXiv:2502.11420, 2025.
\newblock \url{https://arxiv.org/abs/2502.11420}

\bibitem{arxiv:2403.02084}
\newblock \emph{ResAdapter: Domain Consistent Resolution Adapter for Diffusion Models}.
\newblock arXiv:2403.02084, 2024.
\newblock \url{https://arxiv.org/abs/2403.02084}

\bibitem{arxiv:2506.11039}
\newblock \emph{Angle Domain Guidance: Latent Diffusion Requires Rotation Rather Than Extrapolation}.
\newblock arXiv:2506.11039, 2025.
\newblock \url{https://arxiv.org/abs/2506.11039}

\bibitem{arxiv:2407.11942}
\newblock \emph{Context-Guided Diffusion for Out-of-Distribution Molecular and Protein Design}.
\newblock arXiv:2407.11942, 2024.
\newblock \url{https://arxiv.org/abs/2407.11942}

\bibitem{doi:10.3115/981574.981575}
\newblock \emph{Char\_align: A Program for Aligning Parallel Texts at the Character Level}.
\newblock DOI:10.3115/981574.981575, 1993.
\newblock \url{https://doi.org/10.3115/981574.981575}

\bibitem{doi:10.3115/990820.990835}
\newblock \emph{You'll Take the High Road and I'll Take the Low Road: Using a Third Language to Improve Bilingual Word Alignment}.
\newblock DOI:10.3115/990820.990835, 2000.
\newblock \url{https://doi.org/10.3115/990820.990835}

\bibitem{doi:10.3115/1220575.1220585}
\newblock \emph{A Discriminative Matching Approach to Word Alignment}.
\newblock DOI:10.3115/1220575.1220585, 2005.
\newblock \url{https://doi.org/10.3115/1220575.1220585}

\bibitem{arxiv:1705.09655}
\newblock \emph{Style Transfer from Non-Parallel Text by Cross-Alignment}.
\newblock arXiv:1705.09655, 2017.
\newblock \url{https://arxiv.org/abs/1705.09655}

\bibitem{arxiv:2406.05882}
\newblock \emph{Distributional Preference Alignment of LLMs via Optimal Transport}.
\newblock arXiv:2406.05882, 2024.
\newblock \url{https://arxiv.org/abs/2406.05882}

\bibitem{doi:10.1109/tpami.2026.3658965}
\newblock \emph{Diffusion Models and Representation Learning: A Survey}.
\newblock DOI:10.1109/tpami.2026.3658965, 2026.
\newblock \url{https://doi.org/10.1109/tpami.2026.3658965}

\bibitem{arxiv:2505.22165}
\newblock \emph{Unifying Continuous and Discrete Text Diffusion with Non-simultaneous Diffusion Processes}.
\newblock arXiv:2505.22165, 2025.
\newblock \url{https://arxiv.org/abs/2505.22165}

\bibitem{arxiv:2502.08914}
\newblock \emph{Diffusion Models Through a Global Lens: Are They Culturally Inclusive?}.
\newblock arXiv:2502.08914, 2025.
\newblock \url{https://arxiv.org/abs/2502.08914}

\bibitem{arxiv:2503.05149}
\newblock \emph{Development and Enhancement of Text-to-Image Diffusion Models}.
\newblock arXiv:2503.05149, 2025.
\newblock \url{https://arxiv.org/abs/2503.05149}

\end{thebibliography}

\end{document}
